% =====================================================================
%  CL4FMAgents 2026 (NeurIPS 2026 workshop) -- full submission source.
%  Needs in the same Overleaf project: neurips_2026.sty, references.bib, figs/
%  Deadline: 29 Aug 2026. Non-archival. 8 pages excluding references.
% =====================================================================
\documentclass{article}

% Double-blind workshop (anonymous). If CL4FMAgents turns out to be
% SINGLE-blind, swap the two lines below and put the real author block back.
\usepackage[dblblindworkshop]{neurips_2026}
% \usepackage[sglblindworkshop]{neurips_2026}

\usepackage[utf8]{inputenc}
\usepackage[T1]{fontenc}
\usepackage[hidelinks]{hyperref}   % hidelinks = no coloured boxes around links
\usepackage{url}
\usepackage{booktabs}
\usepackage{amsfonts}
\usepackage{amsmath}
\usepackage{amssymb}
\usepackage{nicefrac}
\usepackage{microtype}
\usepackage{xcolor}
\usepackage{graphicx}
\usepackage{subcaption}

\title{Cross-Modal Forgetting and Transfer in Parameter-Efficient\\
Post-Training of a Unified Multimodal Model}

\workshoptitle{NeurIPS 2026 Workshop on Continual Learning in the Era of
Foundation Models and Embodied Agents}

\author{%
  Anonymous Authors \\
}

\begin{document}
\maketitle

\begin{abstract}
Unified multimodal models serve image understanding and image generation from a
single shared backbone, which gives post-training an unusual failure mode:
adapting one capability can degrade another that lives in a \emph{different
modality} of the same weights. We measure this cross-modal forgetting for LoRA
post-training of Janus-Pro-7B. Adapting each direction in turn on five targets
chosen to lie outside the base model's existing competence, from a poetic
register to a language it barely speaks, we find that the capabilities we did not
train stay at base level throughout, on public benchmarks and on free-form
generation alike, even when the adapter drives held-out Welsh perplexity down by
$64\%$. Each adapter changes under $2\%$ of the backbone's weights, concentrated
in the MLP and sparing attention. That localization turns out to cut both ways:
the acquired language is not usable through the visual pathway, leaving Welsh
visual question answering at chance, where the model is not even reading the
question. Adaptation in a unified model is contained, and containment is what
both protects the other modalities and keeps the new capability from reaching
them.
\end{abstract}

\section{Introduction}
Continual learning asks how a model can acquire a new capability without
destroying the ones it already has \citep{mccloskey1989catastrophic,
french1999catastrophic}. Foundation models have made the question practical
rather than academic: post-training is now the normal way to specialize a
released checkpoint, and every such update risks eroding the pretrained
competence that motivated using the checkpoint in the first place
\citep{kirkpatrick2017ewc, li2017lwf}.

Unified multimodal models give this question an unusual shape. Models such as
Chameleon \citep{chameleon2024}, the 4M family \citep{bachmann20244m21},
Unified-IO~2 \citep{lu2024unifiedio2} and Janus \citep{wu2024janus, januspro}
serve both understanding (image$+$text to text) and generation (text to image)
from a single set of weights. When post-training specializes one capability, the
capabilities at risk are not only the neighbouring tasks of the same modality:
they include tasks in a \emph{different} modality that happen to share the same
parameters. We call this cross-modal forgetting. It cannot arise in a task
specialist, such as a text-to-image diffusion model, because there is no other
modality to preserve; it is the natural stress test of the ``one backbone, many
capabilities'' premise.

Parameter-efficient fine-tuning, LoRA in particular \citep{hu2022lora}, is the
default post-training tool and is often assumed to be safe by construction
because it moves few parameters. The assumption is under-tested in exactly the
setting where it matters most. Prior work bounds the problem without answering
it: LoRA forgets less than full fine-tuning while also learning less
\citep{biderman2024lora}, and increasing the LoRA rank beyond a small value buys
little additional task performance \citep{zhang2024scalingmeetsfinetuning}, but
both are measured within a single modality. Work that does address the
understanding/generation tension in unified models resolves it by changing the
model, routing among several adapters or adding a mixture-of-experts layer
\citep{dou2024loramoe}; such designs presuppose that plain LoRA interferes, but
none of them measures how much it actually does. That measurement is what we
provide.

Measuring it credibly requires care about the target. If the fine-tuning target
is something the base model already handles well, the adapter barely has to move
the weights, and finding ``no forgetting'' says nothing: a weak intervention
cannot damage anything, so the result is guaranteed in advance and would be
reported all the same. We therefore select every target by how poorly the base
already handles it, on a scale that is comparable across targets, and we report
how large the resulting weight update is, so that a reader can check that the
intervention was real before believing that nothing broke.

We report two findings on Janus-Pro-7B. First, post-training is strongly
localized: over five targets, LoRA ranks $2$ to $128$ and KL weights
$\lambda\in[0,1.5]$, the capabilities we did not train stay at base level within
evaluation noise, on the four benchmarks, on free-form generation, and in the
other modality's generative distribution. Second, the same localization has a
cost: a language acquired in the text stream is not usable through the visual
pathway at all. Isolation is what protects the untrained capabilities, and it is
also what prevents the new one from reaching them.

\paragraph{Contributions.}
\textbf{(i)} We show that claims of the form ``PEFT does not forget'' are only
meaningful when the adaptation target is demonstrably outside the base model's
existing competence, and we make that condition explicit and measurable: every
target is scored by the base model's perplexity on it relative to a fixed
control, and every adaptation is accompanied by the size of the weight update it
induced.
\textbf{(ii)} Under that condition, the untrained capabilities are preserved
across the whole range of adaptation strength we can reach, including under a
strong distributional shift toward a language the base barely modelled (held-out
perplexity $40.7\to14.6$), and beyond likelihood scoring: free-form English
generation is preserved (MAUVE) and the image-generation distribution is
essentially unmoved (Gen-CMMD $\sim\!1.7\times10^{-3}$). We locate where a cost
first appears: in grounded visual understanding, at the strongest settings.
\textbf{(iii)} A negative transfer result that we argue is the flip side of the
same mechanism: the acquired language improves text-to-image faithfulness only
marginally ($+0.47\pm0.25$ CLIPScore, $3\%$ of the English--Welsh gap) and leaves
Welsh visual question answering at chance, where the model turns out not to be
reading the Welsh question at all.
\textbf{(iv)} Each adapter changes $0.6$ to $1.9\%$ of the backbone's weights,
concentrated in the MLP and sparing attention. The target that needs the smallest
update is not the one furthest from the base: Welsh reaches the largest
perplexity reduction of the three text targets while moving the weights least
($0.6\%$), suggesting the capability is unlocked rather than built.

\section{Related Work}
Continual learning is classically framed as a tension between retaining what a
model can already do and acquiring something new, the stability--plasticity
dilemma \citep{mermillod2013stability}, which remains the organizing question of
the literature on deep networks \citep{delange2022survey, wang2024clsurvey}.
Three families of remedy recur. \emph{Regularization} methods anchor the model to
a reference, in weight space, L2-SP \citep{xuhong2018l2sp} and its
Fisher-weighted form EWC \citep{kirkpatrick2017ewc}, or in function space,
distillation \citep{hinton2015distilling} and Learning-without-Forgetting
\citep{li2017lwf}. \emph{Replay and constraint} methods retain old data and
forbid updates that would hurt it \citep{lopezpaz2017gem, chaudhry2019agem}.
\emph{Parameter-isolation} methods instead give each task its own capacity and
freeze or mask the rest \citep{rusu2016progressive, mallya2018packnet,
serra2018hat}. Our setting inherits from the first and the third: we anchor the
adapter to the frozen base in function space, and LoRA is itself a soft form of
parameter isolation, a small added subspace that leaves the pretrained weights
untouched. That is exactly why the interference it produces is worth measuring
rather than assuming.

What is already known bounds the problem without resolving it here. Forgetting in
large pretrained models is real but scale- and method-dependent: it diminishes
with model scale \citep{ramasesh2022scale}, and parameter-efficient methods learn
less and forget less than full fine-tuning \citep{biderman2024lora}, with rank
beyond a small value buying little \citep{zhang2024scalingmeetsfinetuning}. It has
been documented empirically for language models under continual fine-tuning
\citep{luo2023empirical} and for multimodal ones \citep{zhai2023investigating}.
All of this, however, is measured \emph{within} a modality, or across tasks of the
same kind. A complementary line reduces interference by routing among several
adapters \citep{dou2024loramoe, wu2024mole, li2024mixlora}; those designs
presuppose that plain LoRA interferes, without quantifying how much.

The models that sharpen the question are the unified ones. Fully-tokenized native
models \citep{chameleon2024, bachmann20244m21, lu2024unifiedio2} process all
modalities in one sequence; Janus and Janus-Pro \citep{wu2024janus, januspro}
decouple the visual encoders for understanding and generation while keeping a
single shared backbone. PEFT has been applied to such models
\citep{chern2024anole, serra2024narrowgate, wang2025vora}, but as a means to an
end. To our knowledge, no prior work measures the cross-modal interference of
plain PEFT in a unified multimodal model: what happens to image generation when
the text stream is specialized, and to language, knowledge and visual
understanding when image generation is. Supplying that measurement is the purpose
of this paper. We neither modify the architecture nor propose a new adapter.

\section{Setup}
\label{sec:setup}
\paragraph{Model, directions, and adapter.}
We use Janus-Pro-7B \citep{januspro}: a SigLIP encoder for understanding and a
VQ-VAE (codebook $16{,}384$, $576$ tokens per image) for generation, over one
shared LLaMA-style transformer. It exposes three capability directions:
text-to-text language modelling, image$+$text understanding, and text-to-image
generation. LoRA is injected \emph{only} into the shared backbone
($q,k,v,o,\texttt{gate},\texttt{up},\texttt{down}$), with $\alpha=2r$ fixed
across the rank sweep; the encoders and heads are frozen. The generation forward
traverses only the language-model tensors, while understanding additionally
engages SigLIP.

\paragraph{Controlling how far the adapter drifts: a KL term to the frozen base.}
For both directions we optionally hold the adapter close to the frozen base.
Writing $p_{\text{adapter}}(\cdot\mid x_{<t})$ for the adapted next-token
distribution and $p_{\text{base}}(\cdot\mid x_{<t})$ for the \emph{same} model
with the adapter disabled, the training loss is
\begin{equation}
\mathcal{L}=\underbrace{\mathrm{CE}\big(p_{\text{adapter}},y\big)}_{\text{acquire the target}}
+\lambda\,\frac{1}{T}\sum_{t=1}^{T}\underbrace{\mathrm{KL}\!\big(p_{\text{adapter}}(\cdot\mid x_{<t})\,\big\|\,\mathrm{sg}[p_{\text{base}}(\cdot\mid x_{<t})]\big)}_{\text{stay close to the base}} .
\label{eq:kl}
\end{equation}
The reference is obtained for free by disabling the adapter, so no separate
teacher is stored, and it is held constant by a stop-gradient. The single KL
weight $\lambda$ sets how much the adapter may drift: $\lambda=0$ is plain
fine-tuning, larger $\lambda$ keeps the adapter near the base. We use a
function-space reference because the quantity to preserve is the model's output
\emph{distribution}; an ablation against L2-SP, logit-MSE and the reconstruction
objective of Anole \citep{chern2024anole} at matched update size supports the
choice (Appendix~\ref{app:reg}).

\paragraph{What we evaluate after adaptation.}
We evaluate the directions we did not train with four standard benchmarks scored
by log-likelihood only, hence deterministic and judge-free: WikiText-2 perplexity
\citep{wikitext}, MMLU \citep{mmlu}, POPE \citep{pope} and MME \citep{mme}. The
base POPE accuracy ($0.877$) matches the published Janus-Pro-7B value ($0.874$),
validating the harness. Because likelihood benchmarks never exercise free-form
decoding, we add two distributional generation metrics: MAUVE
\citep{pillutla2021mauve} for text, and, for the image direction under a text
adapter, the CMMD \citep{cmmd} between the adapter's and the base's generations
on a fixed prompt set (Gen-CMMD).

\paragraph{What we mean by retention.}
We call a capability \emph{retained} when the adapter's score on it is within
evaluation noise of the base's. We apply this symmetrically: a large movement is
a behavioural change whatever its sign, so the cases where a score moves
substantially in the \emph{favourable} direction (general-text perplexity, and
MAUVE-vs-base for the medicine adapter) are flagged as changes and analysed
separately in \S\ref{sec:dw}, not counted as evidence of stability. What these
benchmarks can and cannot detect is discussed in \S\ref{sec:boundary} and in the
Limitations.

\paragraph{Choosing targets the base cannot already do.}
As argued in the introduction, certifying that nothing was forgotten is
uninformative if the target was already within the base model's reach. We
therefore score each candidate target by $\times$control: the base model's
perplexity on that corpus divided by its perplexity on a fixed English control.
On this common scale the three text targets are poetry (Project Gutenberg) at
$4.0\times$, Welsh (\texttt{wikimedia/wikipedia}, Welsh subset) at $2.3\times$,
and medical abstracts (\texttt{ccdv/pubmed-summarization}) at $1.2\times$; Welsh
was selected from a wider set of low-resource and constructed languages
(Appendix~\ref{app:langprobe}). The
three therefore differ in \emph{kind}, a register, a language and a domain,
rather than lying on one ordered axis of difficulty, and we will see
(\S\ref{sec:text}) that $\times$control predicts neither how much is acquired nor
how far the weights move. Poetry and medicine were scored in one probe run and
Welsh in another, with the same construction and control corpus. For images we
use a narrow anime style (\texttt{lambdalabs/naruto-blip-captions}, $1{,}141$
pairs) that the base renders weakly, and black-and-white line sketches
(\texttt{zoheb/sketch-scene}). Image acquisition is scored against real style
images with CMMD \citep{cmmd} as the headline, with KID \citep{kid}, DINOv2
\citep{dinov2}, CLIP-IQA \citep{clipiqa} and CLIP directional similarity
\citep{stylegannada} as complements. Full hyperparameters are in
Appendix~\ref{app:impl}. All numbers are single-seed; standard errors quantify
evaluation-set sampling noise at fixed weights (Appendix~\ref{app:impl}).

\section{Results}
We adapt one direction at a time and evaluate every other direction.
Section~\ref{sec:gen} adapts image generation and evaluates text, knowledge and
understanding; Section~\ref{sec:text} adapts language modelling and evaluates
knowledge, understanding and the image-generation distribution. Together the two
form the off-diagonal of a retention matrix over the model's directions.

\begin{figure}[t]
\centering
\begin{subfigure}{0.24\textwidth}
\includegraphics[width=\linewidth]{figs/janus_comparison/no_adapter/human_05_natural.png}
\caption*{\footnotesize Base}
\end{subfigure}
\begin{subfigure}{0.24\textwidth}
\includegraphics[width=\linewidth]{figs/janus_comparison/adapter_solo_gen/human_05_natural.png}
\caption*{\footnotesize Image LoRA}
\end{subfigure}
\begin{subfigure}{0.24\textwidth}
\includegraphics[width=\linewidth]{figs/janus_comparison/no_adapter/human_07_natural.png}
\caption*{\footnotesize Base}
\end{subfigure}
\begin{subfigure}{0.24\textwidth}
\includegraphics[width=\linewidth]{figs/janus_comparison/adapter_solo_gen/human_07_natural.png}
\caption*{\footnotesize Image LoRA}
\end{subfigure}
\caption{The intervention is real: style acquisition on two held-out prompts
(``a young man with blonde spiky hair'', ``an old wizard with a long white
beard''), base versus adapter at identical seeds. Full grids in
Appendix~\ref{app:grids}.}
\label{fig:teaser}
\end{figure}

\subsection{Adapting the generation direction}
\label{sec:gen}
We fine-tune generation to a final configuration ($r=48$, attention and MLP
LoRA, $\lambda=0.5$; the incremental ablation reaching it is in
Appendix~\ref{app:ablation}). The adapter acquires the style
(Figure~\ref{fig:teaser}).

Table~\ref{tab:rank} sweeps the rank from $2$ to $64$, i.e.\ the adapter's
capacity to change the shared backbone. Acquisition improves and saturates
(CMMD $0.139\to0.097$, directional similarity $0.143\to0.158$, knee around rank
$16$ to $32$), while every untrained capability stays at the base level across
the whole range: WikiText perplexity moves within $0.07$ of the base, MMLU, POPE
and MME within sampling noise. More capacity refines acquisition at no measurable
cost to the other modalities. We replicate this on a second, unrelated visual
style (line sketches, Appendix~\ref{app:sketch}); we flag now, rather than leave
to the appendix, that this second style produces the largest footprint of all
five targets and is one of the three settings analysed in \S\ref{sec:boundary}.

\begin{table}[t]
\caption{\textbf{Adapting image generation: rank sweep.} Acquisition (bold =
best) improves and saturates, while every untrained capability stays at the base
level at all ranks. Base POPE ($0.877$) matches the published value ($0.874$).
The base has no directional similarity by construction (DirSim measures the
change away from the base), and its CMMD to the style reference was not computed
in this sweep; Table~\ref{tab:sketch} reports a base CMMD for the second style.
Standard errors are evaluation-set sampling errors and depend on the size of the
evaluation set, not on the adapter, so they are effectively constant along each
row and are given once in the row label rather than in every cell.}
\label{tab:rank}
\centering\small
\begin{tabular}{lccccccc}
\toprule
& Base & $r{=}2$ & $r{=}4$ & $r{=}8$ & $r{=}16$ & $r{=}32$ & $r{=}64$ \\
\midrule
\multicolumn{8}{l}{\emph{Acquisition} (vs.\ real style images)}\\
CMMD $\downarrow$        & n/a & .139 & .124 & .126 & .105 & .108 & \textbf{.097} \\
CLIP DirSim $\uparrow$   & n/a & .143 & .142 & .143 & .154 & .153 & \textbf{.158} \\
\midrule
\multicolumn{8}{l}{\emph{Untrained capabilities}}\\
WikiText PPL $\downarrow$ \tiny($\pm.21$) & 13.71 & 13.66 & 13.70 & 13.76 & 13.74 & 13.78 & 13.76 \\
MMLU $\uparrow$ \tiny($\pm.013$)          & .492 & .491 & .491 & .491 & .501 & .497 & .489 \\
POPE acc $\uparrow$ \tiny($\pm.006$)      & .877 & .879 & .881 & .878 & .878 & .879 & .878 \\
MME $\uparrow$ \tiny($\pm31$)             & 1785 & 1777 & 1777 & 1786 & 1786 & 1796 & 1781 \\
\bottomrule
\end{tabular}
\end{table}

\subsection{Adapting the language-modelling direction}
\label{sec:text}
We now specialize text on three targets that differ in kind: a register
(poetry), a domain (medicine) and a language (Welsh). For each we report
acquisition together with the untrained capabilities and the image-generation
shift.

\paragraph{Sweeping the KL weight.}
Table~\ref{tab:textkl} sweeps $\lambda$ at rank $16$ on poetry. Acquisition is
strongest with little regularization and weakens monotonically as the adapter is
held closer to the base (poetry perplexity $24.5$ at $\lambda=0$ rising to
$26.7$ at $\lambda=1.5$), the expected trade between acquiring and preserving.
Knowledge and POPE stay flat. Crucially, the bottom row shows the \emph{other
modality}: Gen-CMMD stays flat and near zero ($\sim\!1.7\times10^{-3}$) at every
$\lambda$, including $\lambda=0$, where the text objective is entirely
unconstrained. Even with no KL term at all, specializing the text stream does not
measurably move the image-generation distribution. Appendix~\ref{app:t2itext}
shows the same conclusion qualitatively, and Appendix~\ref{app:textsweep} reports
the complementary rank sweep on this target.

\begin{table}[t]
\caption{\textbf{Adapting text (poetry): KL sweep at rank 16.} Acquisition (bold
= best) is strongest at low $\lambda$ and weakens as the KL term tightens. MMLU
and POPE stay flat; MME shows a small dip at $\lambda=0$ (see
\S\ref{sec:boundary}). \emph{Gen-CMMD} is the CMMD between the adapter's and the
base's image generations on a fixed prompt set ($N{=}128$): it stays near zero at
every $\lambda$, so the text adapter does not measurably shift image generation.
Text runs use a subsampled evaluation configuration, so absolute values are
comparable within the table, against the base column, not against
Table~\ref{tab:rank}; the smaller evaluation sets are why the standard errors are
larger here. As in Table~\ref{tab:rank}, standard errors depend on the evaluation
set rather than the adapter and are given once per row.}
\label{tab:textkl}
\centering\small
\begin{tabular}{lcccccc}
\toprule
& Base & $\lambda{=}0$ & $\lambda{=}0.25$ & $\lambda{=}0.5$ & $\lambda{=}1$ & $\lambda{=}1.5$ \\
\midrule
\multicolumn{7}{l}{\emph{Acquisition} (target style)}\\
Poetry PPL $\downarrow$ & 32.7 & \textbf{24.5} & \textbf{24.5} & 24.9 & 25.9 & 26.7 \\
\midrule
\multicolumn{7}{l}{\emph{Untrained capabilities}}\\
WikiText PPL $\downarrow$ \tiny($\pm.21$) & 13.71 & 12.61 & \underline{12.50} & 12.51 & 12.75 & 12.93 \\
MMLU $\uparrow$ \tiny($\pm.021$)          & .481 & .474 & .481 & .479 & .475 & .483 \\
POPE acc $\uparrow$ \tiny($\pm.010$)      & .875 & .868 & .870 & .871 & .878 & .873 \\
MME $\uparrow$ \tiny($\pm42$)             & 1130 & 1101 & 1116 & 1128 & 1127 & 1128 \\
\midrule
\multicolumn{7}{l}{\emph{Shift in the other modality}}\\
Gen-CMMD $\downarrow$ ($\times10^{-3}$) & 0 & 1.73 & 1.71 & 1.73 & 1.77 & 1.65 \\
\bottomrule
\end{tabular}
\end{table}

\paragraph{A domain and a language.}
Poetry and medicine both lie inside the base's competence (poetic English,
medical English), so the absence of forgetting could still be attributed to a
weak intervention. We therefore add a target the base barely speaks: Welsh.
Table~\ref{tab:targets} reports both. Medicine is acquired strongly (held-out
PubMed perplexity $12.8\to8.4$; on the $\times$control scale the domain moves
from $1.2\times$, harder than generic English, to $0.79\times$, easier), and
Welsh more strongly still (held-out perplexity $40.7\to14.6$, $-64\%$); on Welsh
seeds the adapter sustains Welsh where the base slips into English or its
assistant persona (Appendix~\ref{app:welsh}). Qualitative examples for medicine
are in Appendix~\ref{app:medicine}. In both cases the untrained capabilities stay
at base level and Gen-CMMD stays at the same $\sim\!1.7\times10^{-3}$ order as
poetry.

Two observations sharpen the picture. First, \emph{distance from the base
predicts neither outcome}. On the common $\times$control scale Welsh ($2.3\times$)
sits between poetry ($4.0\times$) and medicine ($1.2\times$), yet it produces the
largest perplexity reduction of the three and needs the \emph{smallest} weight
change ($\lVert\Delta W\rVert/\lVert W\rVert=0.6\%$, versus $1.3\%$ for poetry
and $1.7\%$ for medicine), despite the weakest KL term ($\lambda=0.1$). How
unfamiliar a corpus looks to the base model therefore tells us little about how
much of the network has to move to fit it; the new language appears to be
unlocked rather than built, i.e.\ partly latent in the pretrained backbone.
Second, under the adapter, English prompts are still continued in English, so the
language is added without overwriting the existing one.

This is a shift of the output distribution, not competence: the samples
(Appendix~\ref{app:welsh}) sustain Welsh morphology and register where the base
does not, but remain locally repetitive under greedy decoding. We read the
retention result against that weaker premise.

\begin{table}[t]
\caption{\textbf{A distant domain and a new language.} Medicine and Welsh are
each acquired strongly while every untrained capability stays at base level and
image generation is essentially unchanged. The two targets use different
evaluation configurations (full for medicine, subsampled for Welsh), so their
base columns differ in absolute value; each comparison is within its own pair of
columns. Errors are one standard error from evaluation-set resampling. Entries
marked n/a are metrics defined only for the other target.}
\label{tab:targets}
\centering\small
\setlength{\tabcolsep}{5pt}
\begin{tabular}{lcccc}
\toprule
& \multicolumn{2}{c}{Medicine ($r{=}32$, $\lambda{=}0.25$)}
& \multicolumn{2}{c}{Welsh ($r{=}32$, $\lambda{=}0.1$)} \\
\cmidrule(lr){2-3}\cmidrule(lr){4-5}
& Base & LoRA & Base & LoRA \\
\midrule
\multicolumn{5}{l}{\emph{Acquisition}}\\
PubMed PPL $\times$control $\downarrow$ & 1.23$\times$ & \textbf{0.79$\times$} & n/a & n/a \\
Welsh held-out PPL $\downarrow$         & n/a & n/a & 40.7 & \textbf{14.6} \\
$\lVert\Delta W\rVert/\lVert W\rVert$   & --- & 1.7\% & --- & 0.6\% \\
\midrule
\multicolumn{5}{l}{\emph{Untrained capabilities}}\\
WikiText PPL $\downarrow$ & 13.71\,$\pm$.21 & 13.97\,$\pm$.21 & 13.71\,$\pm$.21 & 12.07\,$\pm$.18 \\
MMLU $\uparrow$           & .492\,$\pm$.013 & .495\,$\pm$.013 & .481\,$\pm$.021 & .479\,$\pm$.016 \\
POPE acc $\uparrow$       & .877\,$\pm$.006 & .877\,$\pm$.006 & .875\,$\pm$.010 & .876\,$\pm$.010 \\
MME $\uparrow$            & 1785\,$\pm$31   & 1766\,$\pm$32   & 1130\,$\pm$42   & 1124\,$\pm$42 \\
\midrule
\multicolumn{5}{l}{\emph{Shift in the other modality}}\\
Gen-CMMD $\downarrow$ ($\times10^{-3}$) & 0 & 1.77 & 0 & 1.73 \\
\bottomrule
\end{tabular}
\end{table}

\subsection{Beyond likelihood: free-form generation}
The four benchmarks score multiple-choice and log-likelihood items, never
free-form decoding, so a generative degradation could hide behind flat
likelihood numbers. We close that gap with MAUVE \citep{pillutla2021mauve}, the
text counterpart of our image CMMD (Table~\ref{tab:mauve}). Three readings
follow. \emph{Cross-modal}: the image adapters generate free English essentially
identical to the base (MAUVE-vs-base $0.95$/$0.92$), the generative analogue of
their near-zero Gen-CMMD. \emph{Text, with one caveat}: no adapter's free English
falls below the base against human text, and the text adapters raise it
($0.60$--$0.77$), so none of them degrades English relative to the human
reference. The medicine adapter's MAUVE-vs-base of $0.21$ is nonetheless a large
distributional change: its free English is far from the base's while being closer
to human text. By our definition (\S\ref{sec:setup}) that is a behavioural
change, not retention; what the two columns jointly rule out is degradation.
\emph{Acquisition}: the text adapters move far closer to the target distribution
than the base (poetry $0.45\to0.80$, medical abstracts $0.31\to0.91$, Welsh
$0.14\to0.67$), confirming acquisition at the generation-distribution level and
not only in perplexity. Appendix~\ref{app:textqual} shows the corresponding
poetry continuations.

\begin{table}[t]
\caption{\textbf{Free-form generation (MAUVE)} \citep{pillutla2021mauve}
($[0,1]$, higher = closer; $N\!\le\!2000$, GPT-2-large features). \emph{English
vs.\ base}: image adapters are indistinguishable from the base. \emph{English
vs.\ human}: no adapter falls below the base ($0.38$); text adapters raise it.
\emph{Target acquisition}: text adapters move far closer to the target. Base
continuations are coherent English, not degenerate (Appendix~\ref{app:mauve}).
MAUVE is a divergence between two sets of $N\!\le\!2000$ generations and has no
closed-form standard error; we therefore report it without one and read only
differences larger than the $0.1$ scale on which the entries here separate.}
\label{tab:mauve}
\centering\small
\begin{tabular}{lccc}
\toprule
Adapter & English & English & Target acq. \\
        & vs.\ human & vs.\ base & base\,$\to$\,adapter \\
\midrule
Base (reference)        & 0.38 & --- & --- \\
Anime, image (Final)    & 0.52 & 0.95 & --- \\
Sketch, image           & 0.46 & 0.92 & --- \\
Poetry                  & 0.70 & 0.83 & 0.45\,$\to$\,0.80 \\
Medicine                & 0.60 & 0.21 & 0.31\,$\to$\,0.91 \\
Welsh                   & 0.77 & 0.68 & 0.14\,$\to$\,0.67 \\
\bottomrule
\end{tabular}
\end{table}

\subsection{Where a cost first appears}
\label{sec:boundary}

\paragraph{First, do our measures move at all?}
A result of the form ``nothing changed'' is only worth reading if the
instruments can register change, so before looking for a boundary we check that
they do. Table~\ref{tab:joint} reports a control in which a text stream is
alternated into the image adapter at three corpus sizes. When the text corpus is
small, held-out perplexity inflates dramatically, from $13.71$ at base to
$47.87$ and $55.10$, roughly $60$ standard errors, while the same adapter trained
on a large distant corpus leaves it at $12.46$. The inflation therefore tracks
corpus size, a small corpus being revisited and over-fit, and not domain
distance; the point for us is that the perplexity measure is far from inert, and
that a naive protocol could easily have mistaken this artefact for forgetting.
Knowledge and understanding stay flat even here, which is the first hint that
they are harder to move than perplexity, and also the reason we treat their
flatness elsewhere with more caution than perplexity's.

\begin{table}[t]
\caption{\textbf{Sensitivity control.} A text stream alternated into the image
adapter at three corpus sizes. Held-out perplexity moves by tens of standard
errors when the text corpus is small, so the measure is not inert; the movement
tracks corpus size rather than domain distance. Knowledge and understanding are
unaffected throughout. Full evaluation configuration, single seed.}
\label{tab:joint}
\centering\small
\setlength{\tabcolsep}{6pt}
\begin{tabular}{llcccc}
\toprule
\textbf{Condition} & \textbf{Text corpus} & \textbf{PPL}$\downarrow$ & \textbf{MMLU}$\uparrow$ & \textbf{POPE acc}$\uparrow$ & \textbf{MME}$\uparrow$ \\
 & & \tiny($\pm.21$) & \tiny($\pm.013$) & \tiny($\pm.006$) & \tiny($\pm31$) \\
\midrule
Base (no adapter) & ---                  & $13.71$ & $0.492$ & $0.877$ & $1785$ \\
Final (image-only)& ---                  & $13.80$ & $0.491$ & $0.879$ & $1789$ \\
Joint matched     & near, 320 chunks     & $47.87$ & $0.505$ & $0.876$ & $1747$ \\
Joint far         & distant, 7672 chunks & $12.46$ & $0.491$ & $0.878$ & $1795$ \\
Joint far@320     & distant, 320 chunks  & $55.10$ & $0.481$ & $0.873$ & $1736$ \\
\bottomrule
\end{tabular}
\end{table}

\paragraph{Where the limit is.}
Reporting only preserved capabilities would say nothing about where the limit
is, so we deliberately searched for the settings that break it. Three
independent settings, each at the strongest end of adaptation, show a movement
in the same direction and in the same measure, grounded visual understanding
(MME), while knowledge and language stay flat:
(i) at $\lambda=0$, the unregularized poetry adapter, MME $1130\to1101$
(Table~\ref{tab:textkl}), which the KL term removes;
(ii) pushing poetry to rank $128$, well past the acquisition knee, buys almost
nothing ($23.8\to23.2$) while MME moves $1122\to1105$ and POPE $0.877\to0.867$;
(iii) the hardest visual style, line sketches, is the only target for which the
footprint is uniformly larger than the others, most visibly on MME
($1785\to1713$; Appendix~\ref{app:sketch}).

Each movement lies within roughly one to two standard errors of the difference,
and all runs are single-seed, so none is individually conclusive; their agreement
in sign is equally consistent with a systematic property of MME, which has the
widest dynamic range of our four measures. We report the pattern as an
indication of where to look, not as a finding. What would settle it is set out in
the Limitations.

\subsection{The flip side: what is learned does not cross modalities}
Localized adaptation is usually reported as good news. We show it has a
symmetric cost. Having acquired Welsh in the text direction, can the model
\emph{act} on Welsh in the visual directions? This is a transfer question, not a
retention one: the base is expected to fail on Welsh because it barely knows it,
so any gain must come from the adapter (protocol in Appendix~\ref{app:xmodal}).

Transfer was mechanistically plausible here, which is why the negative answer is
worth reporting. Janus' image-generation forward pass traverses the same
language-model tensors that the adapter modifies, so a text-side update reaches
the generation pathway by construction; nothing in the architecture prevents a
representation acquired in text from conditioning image tokens. A reader who
accepts that a shared backbone carries shared representations should expect at
least partial transfer. We find instead that the pathway is architecturally open
and functionally unused.

Table~\ref{tab:xmodal} shows the answer is essentially no. Generating images from
Welsh prompts, the adapter improves CLIPScore against the same English reference
only marginally over the base ($10.6\to11.0$, $+0.47\pm0.25$, $20/30$ prompts
improved), which closes just $3\%$ of the English--Welsh gap ($15.8$) and stays
far from English-prompt faithfulness. On Welsh visual question answering there is
no transfer at all: on $28$ balanced yes/no items the base scores $53.6\%$ and
the adapter $50.0\%$, both indistinguishable from chance ($95\%$ Wilson intervals
$[36,70]$ and $[33,67]$), against $100\%$ in English for both. The failure is not
a matter of degree. On the Welsh items the base answers ``yes'' to $11$ of the
$14$ questions whose correct answer is ``no'', while answering $12$ of $14$
``yes'' items correctly: it is not reading the Welsh question at all, and falls
back on a response bias that a balanced item set converts into chance-level
accuracy. The English controls are unchanged, so the probe is sound and the
adapter has damaged nothing. Cross-modal grounding of a new language would
require paired image--text data; text-only post-training cannot supply it. The
adapter is localized to the text stream in exactly the sense that protects the
other modalities, and for exactly the same reason it cannot serve them.

\begin{table}[t]
\caption{\textbf{Cross-modal transfer of the acquired language.} On the English
control columns the adapter matches the base, so the probes are sound. On the
Welsh columns the adapter improves text-to-image faithfulness only marginally and
Welsh VQA not at all; both VQA scores are within the $95\%$ Wilson interval of
chance on $n{=}28$ balanced yes/no items. On the Welsh side of (A) the
adapter--base difference is $+0.47\pm0.25$ over $30$ paired prompts. Text-only
acquisition does not reach the visual pathway.}
\label{tab:xmodal}
\centering\small
\begin{tabular}{lcccc}
\toprule
& \multicolumn{2}{c}{Base} & \multicolumn{2}{c}{Welsh LoRA} \\
\cmidrule(lr){2-3}\cmidrule(lr){4-5}
& English & Welsh & English & Welsh \\
\midrule
(A) text$\to$image, CLIPScore $\uparrow$ ($n{=}30$) & 26.4 & 10.6 & 26.6 & 11.0 \\
(B) VQA accuracy $\uparrow$ ($n{=}28$)              & 100\% & 53.6\% & 100\% & 50.0\% \\
\quad\tiny 95\% Wilson interval & \tiny[88,100] & \tiny[36,70] & \tiny[88,100] & \tiny[33,67] \\
\bottomrule
\end{tabular}
\end{table}

\subsection{How large is the intervention?}
\label{sec:dw}
A sceptical reading of everything above is that nothing is forgotten because
nothing is changed. The size of the updates rules that out
(Appendix~\ref{app:dw}). We summarize each update
$\Delta W=(\alpha/r)BA$ by $\lVert\Delta W\rVert_F/\lVert W\rVert_F$, the size of
the update relative to the weights it modifies, a task-vector magnitude
\citep{ilharco2023task}. It is small but real, $0.6$ to $1.9\%$, and it is
\emph{ordered} in both knobs and both directions: rising with rank (image
$0.58\%$ at $r{=}2$ to $1.91\%$ at $r{=}64$; text $0.67$ to $1.66\%$), shrinking
as the KL weight grows (image $1.55$ to $1.33\%$; text $1.15$ to $0.80\%$). These
sub-$2\%$ updates already shift behaviour
substantially (poetry perplexity $32.7\to23.8$; visible style transfer,
Figure~\ref{fig:teaser}) while the other modality stays flat. Per-module heatmaps
(Appendix~\ref{app:dw}) show the change is localized: the attention $q/k/v$
projections are nearly untouched and the update concentrates in the MLP, broadly
over \texttt{gate} for text and toward the final layer's \texttt{down} for image.
The global percentage therefore understates the local effect, and the two
directions differ in \emph{where} they write.

\paragraph{An anomaly we report but do not explain.}
General-text perplexity falls consistently \emph{below} the base after text
adaptation, across ranks, KL weights and all three targets (e.g.\
$13.71\to12.07$ for Welsh), corroborated by the higher MAUVE of the adapters'
free English. We have no confirmed mechanism, and distinguishing the candidates
requires multi-seed runs we have not done, so we report it as an observation and
build no claim on it (Appendix~\ref{app:anomaly}).

\section{Discussion}
\paragraph{Localization is a trade, not a free lunch.}
Read together, our two main results are the same phenomenon seen from two sides.
A small, MLP-localized update in the shared backbone is enough to acquire a
poetic register, a medical domain, or a new language, and it leaves the other
capability directions at base level, on the four benchmarks, in free-form text,
and in the image-generation distribution. That same localization is why the
acquired language cannot be used through the visual pathway. Practitioners who
post-train a unified model can take the first half as reassurance and must take
the second half as a design constraint: if a new capability is to be available in
more than one modality, the evidence here says it has to be trained there,
because it will not arrive by itself.

\paragraph{The KL weight controls drift, not reach.}
The KL term behaves as expected, trading acquisition for preservation
monotonically, and it removes the one movement we could induce ($\lambda=0$,
Table~\ref{tab:textkl}). But it is orthogonal to the transfer question: at every
$\lambda$, including $\lambda=0$, Gen-CMMD is flat, so no setting of it makes the
text update reach the image pathway. How far the model drifts and how far what it
learns travels are controlled by different things.

\paragraph{Limitations.}
Our evidence comes from a single unified model, Janus-Pro-7B, whose
understanding and generation encoders are \emph{decoupled} over a shared
backbone; because LoRA acts only on that backbone, the generation pathway is
partly insulated by construction, which may itself be part of why adaptation
stays localized. Our protocol is also a \emph{single} adaptation step from a
pretrained checkpoint, not a sequence of tasks: we measure interference between
capabilities that coexist in one checkpoint, which is the sub-problem of
continual learning that unified multimodal models make acute, but it is not
sequential continual learning, and we use ``transfer'' in that restricted sense
rather than in the task-sequence sense of \citet{lopezpaz2017gem}. Chaining our
three text targets into one adapter and re-measuring each after every step is the
natural extension and the one we would run next. Whether our conclusions hold for
fully-tokenized early-fusion models (e.g.\ Chameleon), where a single token
stream shares all weights, is untested and is the central open question; we could
not include such models here because their generation is either disabled in the
public release (Chameleon) or inference-only with inaccessible weights (Anole),
the very constraint that motivated our choice of Janus. We do not vary the base
model or its scale.

Our central claim is a negative one, and single-seed evidence cannot establish
it. The standard errors we report throughout are evaluation-set sampling errors
at fixed weights; they do not capture run-to-run variance, which is precisely the
quantity a null result requires. Two experiments would convert what we can only
indicate into a measurement: repeated seeds, and a control in which forgetting is
deliberately induced, showing that these measures move when a capability is
genuinely damaged. Relatedly, our four benchmarks were not selected for
sensitivity to cross-modal damage, because no prior protocol establishes what
such damage looks like in a unified model; the check in Table~\ref{tab:joint}
shows only that perplexity is not inert, and the small movements of
\S\ref{sec:boundary} are suggestive rather than established. Our results should
therefore be read as a failure to detect cross-modal forgetting under this
protocol rather than as a demonstration of its absence. The general-text
improvement is likewise unexplained. Our comparison of anti-forgetting objectives is
scoped to forgetting mitigation; orthogonal adaptation-capacity variants such as
DoRA \citep{liu2024dora} and multi-adapter routing are out of scope. Finally,
free-form quality is assessed distributionally rather than by a human study.

\section{Conclusion}
We measured cross-modal catastrophic forgetting under parameter-efficient
post-training of a unified multimodal model, choosing every adaptation target by
how poorly the base model already handles it so that preservation is certified
against a real intervention. Across two directions, five targets, LoRA ranks $2$
to $128$ and KL weights $\lambda\in[0,1.5]$, the capabilities we did not train
stay at base level within evaluation noise, including under a strong
distributional shift toward a language the base barely modelled, and including in
free-form generation, which likelihood benchmarks do not test. Where a cost
appears at all it is small, confined to grounded visual understanding, and
reachable only at the strongest settings. The same localization that delivers
this stability also confines what is learned: the newly acquired language does
not transfer to the visual pathway, although the architecture leaves that pathway
open. The next steps we consider most valuable are repeated seeds and an
induced-forgetting control, which would turn a failure to detect forgetting into
a measurement of its absence, followed by a sequential version of this protocol
and replication on a fully-tokenized early-fusion model, where the insulation our
conclusions may partly rest on is absent.

\bibliographystyle{plainnat}
\bibliography{references}
\appendix

\section{Implementation details}
\label{app:impl}

\paragraph{Model.}
All experiments use Janus-Pro-7B (\texttt{deepseek-community/Janus-Pro-7B}),
loaded in \texttt{bfloat16} via the Hugging Face
\texttt{JanusForConditionalGeneration} class. Janus-Pro is a unified
autoregressive model with a shared transformer backbone and \emph{decoupled}
modality encoders: a SigLIP vision encoder for understanding (image$\to$text) and
a separate vector-quantized (VQ) tokenizer with an image head for generation
(text$\to$image), where an image is a $24\times24=576$ grid of discrete tokens
from a codebook of size $16384$. We use the two directions through the
processor's text mode (text language modelling) and image mode
(\texttt{max\_new\_tokens}$=576$ then VQ decoding). LoRA is applied \emph{only} to
the shared backbone; the SigLIP and VQ encoders/heads are frozen, which is what
makes the modality-localized weight analysis of Appendix~\ref{app:dw}
meaningful.

\paragraph{LoRA configuration (fixed across all runs).}
Every adapter targets the seven projections of each backbone layer---%
\texttt{q,k,v,o\_proj} (attention) and \texttt{gate,up,down\_proj} (MLP), with
$\alpha=2r$ (so the LoRA scaling $\alpha/r=2$ is constant across the rank sweep),
\texttt{bias=none}, causal-LM task type, and no rsLoRA; the induced update is
$\Delta W=(\alpha/r)BA$. Only the swept variable (rank $r$ or KL weight $\lambda$)
changes within a sweep.

\paragraph{Optimization (common).}
AdamW with $\beta=(0.9,0.95)$ (matching the Janus pre-training $\beta_2$) and zero
weight decay, a cosine schedule with $5\%$ warmup, gradient accumulation $8$,
gradient clipping $1.0$, \texttt{bfloat16}, single H100. With the regularizer the
loss is
$\mathcal{L}=\mathrm{CE}+\lambda\,\mathrm{KL}(p_{\text{adapter}}\|p_{\text{base}})$,
where $p_{\text{base}}$ is a \texttt{disable\_adapter()} no-grad forward on the
same inputs ($\approx$$2\times$ forward cost).

\paragraph{Uncertainty quantification.}
We report each number as $\text{mean}\pm$ one standard error, computed without
retraining and with an estimator-appropriate method per metric type. The accuracy
benchmarks (MMLU, POPE, MME) are proportions of correct items over a fixed
evaluation set, so their standard error follows the binomial/Wilson form
\citep{wilson1927,brown2001interval}. Per-instance metrics averaged over the
generated set (CLIP directional similarity, DINO similarity, CLIP-IQA) use
$s/\sqrt{n}$ over the $n$ generations, equivalently a nonparametric bootstrap
\citep{efron1993bootstrap}. The set-level distributional metrics (CMMD, KID) have
no closed-form variance, so their standard error is the standard deviation of a
nonparametric bootstrap over the generated and reference sets ($B{=}2000$
resamples) \citep{efron1993bootstrap,binkowski2018demystifying}. Perplexity is a
monotone transform of the mean per-chunk negative log-likelihood, whose standard
error propagates through the exponential. All standard errors quantify
evaluation-set sampling noise at fixed model weights; they do not capture
training-seed variance. The small benchmark movements we observe (e.g.\ MMLU
$0.492\!\to\!0.479$, within $\sim\!1$ standard error) thus lie within sampling
noise, whereas the acquisition shifts (e.g.\ CMMD, DirSim) do not. This is the
basis for the cautious reading of the boundary settings in
Section~\ref{sec:boundary}.

\paragraph{Text (poetry) runs.}
Data: \texttt{biglam/gutenberg-poetry-corpus}, grouped by source document with
front-matter removed and windowed to $512$ tokens ($6000$ train / $200$ val
windows). Fixed: $r{=}16$, $\alpha{=}32$, dropout $0.05$, LR $10^{-4}$, $3$
epochs ($150$-min wall-clock budget). Rank sweep: $r\in\{4,8,16,32,64\}$ at
$\lambda{=}0.5$; $\lambda$ sweep: $\lambda\in\{0,0.25,0.5,1,1.5\}$ at $r{=}16$.

\paragraph{Image (Naruto) runs.}
Data: \texttt{lambdalabs/naruto-blip-captions} ($\approx$$1141$ image--caption
pairs); acquisition on $80$ held-out prompts. Final configuration
(Table~\ref{tab:panel}): $r{=}48$, $\alpha{=}96$, dropout $0.10$, LR $10^{-4}$,
$6$ epochs, no caption prefix, $\lambda{=}0.5$. Rank sweep:
$r\in\{2,\dots,64\}$ at $\lambda{=}0.5$; $\lambda$ sweep at $r{=}32$.

\paragraph{Medicine, Welsh and sketch runs.}
Medicine: \texttt{ccdv/pubmed-summarization}, $r{=}32$, $\alpha{=}64$,
$\lambda{=}0.25$. Welsh: \texttt{wikimedia/wikipedia} (Welsh subset), $r{=}32$,
$\alpha{=}64$, $\lambda{=}0.1$, $3$ epochs. Sketch: \texttt{zoheb/sketch-scene},
identical configuration to the Naruto Final run ($r{=}48$, $\alpha{=}96$,
$\lambda{=}0.5$, LR $10^{-4}$, $6$ epochs). All three share the common
optimization settings above.
% TODO FC: add LR / epochs / corpus size for medicine, and MAX_DOCS /
% TRAIN_WINDOWS for Welsh, from the SLURM run scripts.

\paragraph{Prompts and image generation.}
Janus generates images autoregressively from a natural-language prompt
(text$\to$image); every generated image in this paper is therefore conditioned on
an explicit text prompt. During training the conditioning text is the dataset's
BLIP caption; the \textit{Capacity} and \textit{+KL} configurations prepend a
fixed style trigger (``anime sprite in Naruto style, ''), which \textit{Scaled}
and \textit{Final} remove, so the Final adapter applies the style
\emph{unconditionally}, with no trigger token at inference. Acquisition metrics
(Table~\ref{tab:panel}) are computed on images generated from $80$ held-out Naruto
captions. The qualitative grids instead use a fixed, curated set of generic
natural-language prompts mixing in-domain and clearly out-of-domain subjects
(e.g.\ ``a young man with blonde spiky hair'', ``a samurai warrior holding a
katana'', ``an old wizard with a long white beard''), chosen to test whether the
style transfers to arbitrary content; the exact prompt is printed beneath each row
of every grid (Figures~\ref{fig:janus_qualitative_comparison},
\ref{fig:rank_visual}, \ref{fig:kl_visual}). The cross-modal check under the text
adapter (Figure~\ref{fig:t2i_text_kl}) uses four neutral prompts (e.g.\ ``a cat
sitting on a chair''). All conditions share the same seed schedule, so within a
row the adapter is the only source of variation.


\section{Selecting the language target}
\label{app:langprobe}
Before training anything, we scored candidate language targets by
$\times$control: the base model's perplexity on the candidate corpus divided by
its perplexity on an English WikiText control (control perplexity $14.22$ under
this probe's chunking). Candidates were real, available corpora: low-resource
natural languages via Wikipedia, and a synthetic substitution-cipher language
built from English as a maximally distant anchor. Table~\ref{tab:langprobe}
reports the result.

We did not take the most distant candidate. The cipher is by far the furthest
($10.3\times$), but it is a deterministic character-level remapping of English:
an adapter can fit it by learning the substitution table, which would make the
resulting ``acquisition'' a shortcut rather than a new linguistic capability.
Klingon is marginally further than Welsh ($2.51\times$ vs.\ $2.34\times$) but its
available corpus is a few thousand short sentences, too small to separate
acquisition from memorization at our training budget. Welsh combines a genuine
distance from the base with an abundant, naturally occurring corpus, which is why
we selected it.

\begin{table}[h]\centering\small
\setlength{\tabcolsep}{8pt}
\begin{tabular}{lccl}
\toprule
Candidate & Base PPL & $\times$control & Note \\
\midrule
Cipher (substitution) & $146.8$ & $10.33$ & deterministic remapping of English \\
Klingon               & $35.7$  & $2.51$  & constructed; corpus too small \\
\textbf{Welsh}        & $33.3$  & $\mathbf{2.34}$ & \textbf{selected}: distant and abundant \\
Yoruba                & $24.0$  & $1.69$  & closer to the base \\
\bottomrule
\end{tabular}
\caption{Language-target probe, base model only, no training. $\times$control is
the base perplexity on the candidate divided by the base perplexity on an English
control ($14.22$ here), the same construction used for the poetry and medicine
targets. Higher = the base is more surprised by the corpus. We selected Welsh
rather than the two more distant candidates for the reasons given in the text.}
\label{tab:langprobe}
\end{table}


\section{Image-generation configuration: incremental ablation}
\label{app:ablation}
We reach the Final configuration by adding components one or a few at a time from
a naive starting point (attention-only LoRA, rank 8, with a fixed caption-prefix
trigger), which degraded the base and is omitted from Table~\ref{tab:panel}.
\textit{Capacity} restores trainable capacity (rank 16, LoRA on attention and
MLP); \textit{+KL} adds the KL-to-base regularizer ($\lambda{=}0.5$);
\textit{Scaled} removes the training-time trigger prefix and raises rank, learning
rate and dropout (rank 32, $8\times10^{-5}$, dropout 0.1, six epochs);
\textit{Final} raises the rank to 48 and the learning rate to $10^{-4}$. The
\textit{Joint} conditions additionally alternate a $50/50$ domain-controlled text
stream (matched: Naruto franchise; far: physics and mathematics) into the same
adapter. Table~\ref{tab:panel} reports the acquisition panel; the corresponding
generations are in Figure~\ref{fig:janus_qualitative_comparison}.

\begin{table}[t]
\centering\small
\setlength{\tabcolsep}{5pt}
\renewcommand{\arraystretch}{1.05}
\begin{tabular}{lcccccc}
\toprule
 & \textbf{Capacity} & \textbf{+KL} & \textbf{Scaled} & \textbf{Final} & \textbf{Joint} & \textbf{Joint} \\
 & & & & & \textbf{(matched)} & \textbf{(far)} \\
\midrule
\multicolumn{7}{l}{\textit{Configuration (bold = changed vs.\ previous column)}}\\
Rank $r$              & 16 & 16 & \textbf{32} & \textbf{48} & 48 & 48 \\
LR ($\times10^{-5}$)  & 5 & 5 & \textbf{8} & \textbf{10} & 10 & 10 \\
Epochs               & 8 & 8 & \textbf{6} & 6 & 6 & 6 \\
LoRA dropout         & .05 & .05 & \textbf{.10} & .10 & .10 & .10 \\
Train prefix         & yes & yes & \textbf{no} & no & no & no \\
KL $\lambda$         & 0 & \textbf{0.5} & 0.5 & 0.5 & 0.5 & 0.5 \\
Joint text           & --- & --- & --- & --- & \textbf{matched} & \textbf{far} \\
\midrule
\multicolumn{7}{l}{\textit{Acquisition (vs.\ real style, $N{=}80$; bold = best)}}\\
KID $\downarrow$       & $.039$ & $.017$ & $.016$ & $.016$ & $.016$ & $\mathbf{.014}$ \\
CMMD $\downarrow$      & $.189$ & $\mathbf{.105}$ & $.120$ & $.110$ & $.131$ & $.134$ \\
DINO sim $\uparrow$    & $.629$ & $\mathbf{.663}$ & $.624$ & $.640$ & $.633$ & $.637$ \\
CLIP-IQA $\uparrow$    & $.659$ & $\mathbf{.750}$ & $.694$ & $.701$ & $.725$ & $.704$ \\
CLIP DirSim $\uparrow$ & $-.059$ & $-.033$ & $.157$ & $\mathbf{.163}$ & $.146$ & $.151$ \\
\bottomrule
\end{tabular}
\caption{Image-generation configuration as an incremental ablation: each column
adds one or more components to the previous one to reach \textit{Final}.
\textit{Capacity}: rank 16 with LoRA on attention and MLP. \textit{+KL}: adds the
KL-to-base regularizer. \textit{Scaled}: removes the training-time caption prefix
and raises rank, learning rate and dropout. \textit{Final}: rank 48, LR
$10^{-4}$. \textit{Joint}: alternates a $50/50$ domain-controlled text stream
(matched: Naruto franchise; far: physics/mathematics) into the same adapter. The
naive start (attention-only, rank 8, prefix-gated) degraded the base and is
omitted; CLIP DirSim turning positive at \textit{Scaled} marks the onset of
genuine style transfer.}
\label{tab:panel}
\end{table}

\paragraph{The anchor strength has an acquisition optimum.}
Table~\ref{tab:klimg} sweeps $\lambda$ at rank 32 in the image direction. Unlike
the text direction, acquisition here is a single-peaked function of $\lambda$: it
peaks near $\lambda{=}0.25$ to $0.5$ and collapses for $\lambda\!\geq\!1$, where
the adapter is over-constrained toward the base. $\lambda{=}0$ is already
workable, so in this direction the anchor refines rather than enables the recipe.
Retention stays at the base level across the whole $\lambda$ range, so the anchor
changes acquisition without any measurable cross-capability cost. The CLIP style
metrics place the optimum at $\lambda{=}0.25$, yet we adopt $\lambda{=}0.5$ for
the Final recipe because those metrics reward proximity to the style
\emph{distribution} but not whether the prompt's content survives the
stylisation; Appendix~\ref{app:metricdebug} settles this with two content-aware
embedding measures and an independent VLM judge.

\begin{table}[t]
\centering\small
\setlength{\tabcolsep}{6pt}
\begin{tabular}{lcccccc}
\toprule
Metric & Base & $\lambda{=}0$ & $\lambda{=}0.25$ & $\lambda{=}0.5$ & $\lambda{=}1.0$ & $\lambda{=}1.5$ \\
\midrule
\multicolumn{7}{l}{\textit{Acquisition (vs.\ real style images)}}\\
CMMD $\downarrow$      & --- & $.155$ & $\mathbf{.100}$ & $.108$ & $.205$ & $.221$ \\
DINO sim $\uparrow$    & --- & $.626$ & $\mathbf{.647}$ & $.643$ & $.439$ & $.388$ \\
CLIP DirSim $\uparrow$ & --- & $.157$ & $\mathbf{.164}$ & $.152$ & $.049$ & $.024$ \\
CLIP-IQA $\uparrow$    & --- & $.666$ & $.691$ & $.691$ & $.738$ & $\mathbf{.739}$ \\
\midrule
\multicolumn{7}{l}{\textit{Untrained capabilities}}\\
WikiText PPL $\downarrow$ \tiny($\pm.21$) & $13.71$ & $13.67$ & $13.68$ & $13.78$ & $13.76$ & $13.76$ \\
MMLU $\uparrow$ \tiny($\pm.013$)          & $.492$ & $.495$ & $.493$ & $.497$ & $.493$ & $.492$ \\
POPE acc $\uparrow$ \tiny($\pm.006$)      & $.877$ & $.879$ & $.877$ & $.879$ & $.877$ & $.873$ \\
MME $\uparrow$ \tiny($\pm31$)             & $1785$ & $1783$ & $1776$ & $1796$ & $1802$ & $1810$ \\
\bottomrule
\end{tabular}
\caption{Image-generation $\lambda$ sweep at rank 32. \textit{Acquisition} (bold
= best) peaks near $\lambda{=}0.25$ to $0.5$ and degrades for
$\lambda\!\geq\!1$. \textit{Retention} stays at the base level across the whole
$\lambda$ range (same probe configuration as Table~\ref{tab:rank}), so the anchor
changes acquisition without any measurable cross-capability cost.}
\label{tab:klimg}
\end{table}

\begin{figure*}[p]
\centering
\setlength{\tabcolsep}{1.5pt}
\renewcommand{\arraystretch}{1.0}
\small
\begin{tabular}{ccccccc}
\textbf{No adapter} & \textbf{Capacity} & \textbf{+KL} & \textbf{Scaled} & \textbf{Final} & \textbf{Joint (matched)} & \textbf{Joint (far)} \\
\textit{base 7B} & \textit{r16, +MLP} & \textit{+ KL reg} & \textit{r32, untriggered} & \textit{r48, LR $10^{-4}$} & \textit{+ Naruto wiki} & \textit{+ physics wiki} \\
\midrule
\includegraphics[width=0.135\textwidth]{figs/janus_comparison/no_adapter/eval_00_natural.png} &
\includegraphics[width=0.135\textwidth]{figs/janus_comparison/adapter_v2/eval_00_natural.png} &
\includegraphics[width=0.135\textwidth]{figs/janus_comparison/adapter_v3/eval_00_natural.png} &
\includegraphics[width=0.135\textwidth]{figs/janus_comparison/adapter_v4/eval_00_natural.png} &
\includegraphics[width=0.135\textwidth]{figs/janus_comparison/adapter_solo_gen/eval_00_natural.png} &
\includegraphics[width=0.135\textwidth]{figs/janus_comparison/adapter_text_matched/eval_00_natural.png} &
\includegraphics[width=0.135\textwidth]{figs/janus_comparison/adapter_text_far/eval_00_natural.png} \\
\multicolumn{7}{l}{\footnotesize\textit{Prompt (eval\_00):} a red panda with a black face and brown eyes} \\
\addlinespace[2pt]
\includegraphics[width=0.135\textwidth]{figs/janus_comparison/no_adapter/eval_01_natural.png} &
\includegraphics[width=0.135\textwidth]{figs/janus_comparison/adapter_v2/eval_01_natural.png} &
\includegraphics[width=0.135\textwidth]{figs/janus_comparison/adapter_v3/eval_01_natural.png} &
\includegraphics[width=0.135\textwidth]{figs/janus_comparison/adapter_v4/eval_01_natural.png} &
\includegraphics[width=0.135\textwidth]{figs/janus_comparison/adapter_solo_gen/eval_01_natural.png} &
\includegraphics[width=0.135\textwidth]{figs/janus_comparison/adapter_text_matched/eval_01_natural.png} &
\includegraphics[width=0.135\textwidth]{figs/janus_comparison/adapter_text_far/eval_01_natural.png} \\
\multicolumn{7}{l}{\footnotesize\textit{Prompt (eval\_01):} a man with glasses and a hoodie} \\
\addlinespace[2pt]
\includegraphics[width=0.135\textwidth]{figs/janus_comparison/no_adapter/eval_02_natural.png} &
\includegraphics[width=0.135\textwidth]{figs/janus_comparison/adapter_v2/eval_02_natural.png} &
\includegraphics[width=0.135\textwidth]{figs/janus_comparison/adapter_v3/eval_02_natural.png} &
\includegraphics[width=0.135\textwidth]{figs/janus_comparison/adapter_v4/eval_02_natural.png} &
\includegraphics[width=0.135\textwidth]{figs/janus_comparison/adapter_solo_gen/eval_02_natural.png} &
\includegraphics[width=0.135\textwidth]{figs/janus_comparison/adapter_text_matched/eval_02_natural.png} &
\includegraphics[width=0.135\textwidth]{figs/janus_comparison/adapter_text_far/eval_02_natural.png} \\
\multicolumn{7}{l}{\footnotesize\textit{Prompt (eval\_02):} a person with black hair and a green shirt} \\
\addlinespace[2pt]
\includegraphics[width=0.135\textwidth]{figs/janus_comparison/no_adapter/human_05_natural.png} &
\includegraphics[width=0.135\textwidth]{figs/janus_comparison/adapter_v2/human_05_natural.png} &
\includegraphics[width=0.135\textwidth]{figs/janus_comparison/adapter_v3/human_05_natural.png} &
\includegraphics[width=0.135\textwidth]{figs/janus_comparison/adapter_v4/human_05_natural.png} &
\includegraphics[width=0.135\textwidth]{figs/janus_comparison/adapter_solo_gen/human_05_natural.png} &
\includegraphics[width=0.135\textwidth]{figs/janus_comparison/adapter_text_matched/human_05_natural.png} &
\includegraphics[width=0.135\textwidth]{figs/janus_comparison/adapter_text_far/human_05_natural.png} \\
\multicolumn{7}{l}{\footnotesize\textit{Prompt (human\_05):} a young man with blonde spiky hair} \\
\addlinespace[2pt]
\includegraphics[width=0.135\textwidth]{figs/janus_comparison/no_adapter/human_06_natural.png} &
\includegraphics[width=0.135\textwidth]{figs/janus_comparison/adapter_v2/human_06_natural.png} &
\includegraphics[width=0.135\textwidth]{figs/janus_comparison/adapter_v3/human_06_natural.png} &
\includegraphics[width=0.135\textwidth]{figs/janus_comparison/adapter_v4/human_06_natural.png} &
\includegraphics[width=0.135\textwidth]{figs/janus_comparison/adapter_solo_gen/human_06_natural.png} &
\includegraphics[width=0.135\textwidth]{figs/janus_comparison/adapter_text_matched/human_06_natural.png} &
\includegraphics[width=0.135\textwidth]{figs/janus_comparison/adapter_text_far/human_06_natural.png} \\
\multicolumn{7}{l}{\footnotesize\textit{Prompt (human\_06):} a samurai warrior holding a katana} \\
\addlinespace[2pt]
\includegraphics[width=0.135\textwidth]{figs/janus_comparison/no_adapter/human_07_natural.png} &
\includegraphics[width=0.135\textwidth]{figs/janus_comparison/adapter_v2/human_07_natural.png} &
\includegraphics[width=0.135\textwidth]{figs/janus_comparison/adapter_v3/human_07_natural.png} &
\includegraphics[width=0.135\textwidth]{figs/janus_comparison/adapter_v4/human_07_natural.png} &
\includegraphics[width=0.135\textwidth]{figs/janus_comparison/adapter_solo_gen/human_07_natural.png} &
\includegraphics[width=0.135\textwidth]{figs/janus_comparison/adapter_text_matched/human_07_natural.png} &
\includegraphics[width=0.135\textwidth]{figs/janus_comparison/adapter_text_far/human_07_natural.png} \\
\multicolumn{7}{l}{\footnotesize\textit{Prompt (human\_07):} an old wizard with a long white beard} \\
\addlinespace[2pt]
\includegraphics[width=0.135\textwidth]{figs/janus_comparison/no_adapter/human_08_natural.png} &
\includegraphics[width=0.135\textwidth]{figs/janus_comparison/adapter_v2/human_08_natural.png} &
\includegraphics[width=0.135\textwidth]{figs/janus_comparison/adapter_v3/human_08_natural.png} &
\includegraphics[width=0.135\textwidth]{figs/janus_comparison/adapter_v4/human_08_natural.png} &
\includegraphics[width=0.135\textwidth]{figs/janus_comparison/adapter_solo_gen/human_08_natural.png} &
\includegraphics[width=0.135\textwidth]{figs/janus_comparison/adapter_text_matched/human_08_natural.png} &
\includegraphics[width=0.135\textwidth]{figs/janus_comparison/adapter_text_far/human_08_natural.png} \\
\multicolumn{7}{l}{\footnotesize\textit{Prompt (human\_08):} a teenage girl with pink hair} \\
\end{tabular}
\caption{Qualitative comparison across the seven training conditions on
Janus-Pro-7B. Columns 1--5 trace the image-only LoRA progression to the
\textit{Final} configuration; columns 6--7 add a joint text-training stream from
two semantically distant Wikipedia domains (Naruto franchise vs.\
physics/mathematics). All inferences share the same seed schedule; the adapter is
the only source of variation within a row.}
\label{fig:janus_qualitative_comparison}
\end{figure*}

\paragraph{Joint training: a corpus-size effect on perplexity.}
The \textit{Joint} conditions are the sensitivity control promoted to the main
text (Table~\ref{tab:joint}). Adding a text stream inflates held-out text
perplexity only when the text corpus is small: the controlled far@320 run
(distant domain, matched chunk budget) reaches $55.1$, higher than the
near-domain matched run, while MMLU and POPE stay flat. The inflation is
therefore a corpus-size effect, a small corpus being revisited and over-fit, and
not a domain-distance effect.


\section{Generalisation to a second visual style}
\label{app:sketch}
To test that the generation-direction result is not specific to one style, we
replicate the image protocol on a second, unrelated dataset: black-and-white
line-sketch scenes (\texttt{zoheb/sketch-scene}), using the \emph{same}
configuration as Naruto ($r{=}48$, $\alpha{=}96$, $\lambda{=}0.5$, LR
$10^{-4}$). The core result holds (Table~\ref{tab:sketch}): the adapter acquires
the style at the distribution level (CMMD $0.253\!\to\!0.159$, KID
$0.263\!\to\!0.096$, DINO similarity to the sketch reference
$0.133\!\to\!0.371$), while the language and knowledge probes show no
catastrophic forgetting (WikiText PPL $13.71\!\to\!14.18$, MMLU
$0.492\!\to\!0.479$).

This replication also surfaces two findings we use in the main text. First,
\emph{the PEFT configuration is dataset-dependent}: the recipe that saturates a
high-frequency, colour-saturated style only partially imprints the sparse,
low-frequency sketch style (Figure~\ref{fig:sketch}). The per-prompt edit
strength is about half (CLIP DirSim $0.086$, versus $\approx\!0.15$ for Naruto at
the same $\lambda$), and it carries a measurable image-quality cost (CLIP-IQA
$0.747\!\to\!0.592$). Scaling the recipe does not help: a stronger configuration
(rank $64$, $\lambda{=}0.25$) moves the aggregate distribution closer to the
reference (KID $0.096\!\to\!0.065$, DINO $0.371\!\to\!0.521$) but does not
improve the per-prompt edit strength (DirSim $0.086\!\to\!0.089$) and lowers
image quality (CLIP-IQA $0.59\!\to\!0.49$), indicating the limit is intrinsic to
the low-frequency style rather than a matter of capacity. Second, this is the
target whose retention footprint is uniformly the largest of all five, and most
visible on MME ($1785\!\to\!1713$, versus flat for Naruto): it is the third of
the three convergent settings discussed in Section~\ref{sec:boundary}.

\begin{table}[t]
\centering\small
\setlength{\tabcolsep}{8pt}
\begin{tabular}{lcc}
\toprule
Metric & Base & Sketch LoRA \\
\midrule
\multicolumn{3}{l}{\textit{Acquisition (vs.\ sketch reference, $n{=}80$)}}\\
CMMD $\downarrow$      & $.253$ & $\mathbf{.159}_{\pm.007}$ \\
KID $\downarrow$       & $.263$ & $\mathbf{.096}_{\pm.013}$ \\
DINO sim $\uparrow$    & $.133$ & $\mathbf{.371}_{\pm.021}$ \\
CLIP DirSim $\uparrow$ & ---    & $.086_{\pm.006}$ \\
\midrule
\multicolumn{3}{l}{\textit{Retention}}\\
WikiText PPL $\downarrow$ & $13.71_{\pm.21}$ & $14.18_{\pm.22}$ \\
MMLU $\uparrow$           & $.492_{\pm.013}$ & $.479_{\pm.013}$ \\
POPE acc $\uparrow$       & $.877_{\pm.006}$ & $.871_{\pm.006}$ \\
MME $\uparrow$            & $1785_{\pm31}$   & $1713_{\pm31}$ \\
\bottomrule
\end{tabular}
\caption{Generalisation to a second image dataset (line-sketch scenes,
\texttt{zoheb/sketch-scene}) with the \emph{same} configuration as Naruto. The
style is acquired at the distribution level and the language and knowledge probes
show no catastrophic forgetting; the per-prompt edit strength (CLIP DirSim
$0.086$) is about half of Naruto's, evidence that the PEFT recipe is
dataset-dependent. This is also the largest retention footprint of the five
targets, most visible on MME. FID/ArtFID are omitted (numerically unstable at
$n{=}80$).}
\label{tab:sketch}
\end{table}

\begin{figure}[t]
\centering
\begin{subfigure}{0.24\textwidth}
\includegraphics[width=\linewidth]{figs/sketch_comparison/zebra_before.png}
\caption*{\footnotesize Base}
\end{subfigure}
\begin{subfigure}{0.24\textwidth}
\includegraphics[width=\linewidth]{figs/sketch_comparison/zebra_after.png}
\caption*{\footnotesize Sketch LoRA}
\end{subfigure}
\begin{subfigure}{0.24\textwidth}
\includegraphics[width=\linewidth]{figs/sketch_comparison/kite_before.png}
\caption*{\footnotesize Base}
\end{subfigure}
\begin{subfigure}{0.24\textwidth}
\includegraphics[width=\linewidth]{figs/sketch_comparison/kite_after.png}
\caption*{\footnotesize Sketch LoRA}
\end{subfigure}
\caption{Generalisation to line-sketch scenes (same prompt and seed, base vs.\
adapter). The style transfers for some prompts (left pair: white canvas, ink
linework) but only partially for others (right pair: sketch strokes appear yet
photographic colour persists), consistent with the weaker, dataset-dependent
acquisition in Table~\ref{tab:sketch}.}
\label{fig:sketch}
\end{figure}


\section{Qualitative grids}
\label{app:grids}

\begin{figure*}[p]
\centering
\setlength{\tabcolsep}{1.5pt}\small
\renewcommand{\arraystretch}{1.0}
\begin{tabular}{ccccccc}
\textbf{Base} & \textbf{r=2} & \textbf{r=4} & \textbf{r=8} & \textbf{r=16} & \textbf{r=32} & \textbf{r=64} \\
\midrule
\includegraphics[width=0.135\textwidth]{figs/rank_comparison/base/eval_00_natural.png} &
\includegraphics[width=0.135\textwidth]{figs/rank_comparison/r2/eval_00_natural.png} &
\includegraphics[width=0.135\textwidth]{figs/rank_comparison/r4/eval_00_natural.png} &
\includegraphics[width=0.135\textwidth]{figs/rank_comparison/r8/eval_00_natural.png} &
\includegraphics[width=0.135\textwidth]{figs/rank_comparison/r16/eval_00_natural.png} &
\includegraphics[width=0.135\textwidth]{figs/rank_comparison/r32/eval_00_natural.png} &
\includegraphics[width=0.135\textwidth]{figs/rank_comparison/r64/eval_00_natural.png} \\
\multicolumn{7}{l}{\footnotesize\textit{Prompt:} a red panda with a black face and brown eyes} \\
\addlinespace[2pt]
\includegraphics[width=0.135\textwidth]{figs/rank_comparison/base/eval_01_natural.png} &
\includegraphics[width=0.135\textwidth]{figs/rank_comparison/r2/eval_01_natural.png} &
\includegraphics[width=0.135\textwidth]{figs/rank_comparison/r4/eval_01_natural.png} &
\includegraphics[width=0.135\textwidth]{figs/rank_comparison/r8/eval_01_natural.png} &
\includegraphics[width=0.135\textwidth]{figs/rank_comparison/r16/eval_01_natural.png} &
\includegraphics[width=0.135\textwidth]{figs/rank_comparison/r32/eval_01_natural.png} &
\includegraphics[width=0.135\textwidth]{figs/rank_comparison/r64/eval_01_natural.png} \\
\multicolumn{7}{l}{\footnotesize\textit{Prompt:} a man with glasses and a hoodie} \\
\addlinespace[2pt]
\includegraphics[width=0.135\textwidth]{figs/rank_comparison/base/human_05_natural.png} &
\includegraphics[width=0.135\textwidth]{figs/rank_comparison/r2/human_05_natural.png} &
\includegraphics[width=0.135\textwidth]{figs/rank_comparison/r4/human_05_natural.png} &
\includegraphics[width=0.135\textwidth]{figs/rank_comparison/r8/human_05_natural.png} &
\includegraphics[width=0.135\textwidth]{figs/rank_comparison/r16/human_05_natural.png} &
\includegraphics[width=0.135\textwidth]{figs/rank_comparison/r32/human_05_natural.png} &
\includegraphics[width=0.135\textwidth]{figs/rank_comparison/r64/human_05_natural.png} \\
\multicolumn{7}{l}{\footnotesize\textit{Prompt:} a young man with blonde spiky hair} \\
\addlinespace[2pt]
\includegraphics[width=0.135\textwidth]{figs/rank_comparison/base/human_07_natural.png} &
\includegraphics[width=0.135\textwidth]{figs/rank_comparison/r2/human_07_natural.png} &
\includegraphics[width=0.135\textwidth]{figs/rank_comparison/r4/human_07_natural.png} &
\includegraphics[width=0.135\textwidth]{figs/rank_comparison/r8/human_07_natural.png} &
\includegraphics[width=0.135\textwidth]{figs/rank_comparison/r16/human_07_natural.png} &
\includegraphics[width=0.135\textwidth]{figs/rank_comparison/r32/human_07_natural.png} &
\includegraphics[width=0.135\textwidth]{figs/rank_comparison/r64/human_07_natural.png} \\
\multicolumn{7}{l}{\footnotesize\textit{Prompt:} an old wizard with a long white beard} \\
\addlinespace[2pt]
\includegraphics[width=0.135\textwidth]{figs/rank_comparison/base/human_08_natural.png} &
\includegraphics[width=0.135\textwidth]{figs/rank_comparison/r2/human_08_natural.png} &
\includegraphics[width=0.135\textwidth]{figs/rank_comparison/r4/human_08_natural.png} &
\includegraphics[width=0.135\textwidth]{figs/rank_comparison/r8/human_08_natural.png} &
\includegraphics[width=0.135\textwidth]{figs/rank_comparison/r16/human_08_natural.png} &
\includegraphics[width=0.135\textwidth]{figs/rank_comparison/r32/human_08_natural.png} &
\includegraphics[width=0.135\textwidth]{figs/rank_comparison/r64/human_08_natural.png} \\
\multicolumn{7}{l}{\footnotesize\textit{Prompt:} a teenage girl with pink hair} \\
\end{tabular}
\caption{Qualitative rank sweep (Final configuration): the target style is already
present at $r{=}2$ and only refines with rank, remaining stable across the range.}
\label{fig:rank_visual}
\end{figure*}

\begin{figure}[t]
\centering
\includegraphics[width=\textwidth]{figs/rank_sweep/rank_sweep_figure.pdf}
\caption{Rank sweep on the Final configuration (absolute metric values).
\textbf{Left:} visual-style acquisition (CLIP directional similarity); the
dotted line at $0$ marks an adapter with no directional effect. Acquisition is
positive at every rank and rises modestly, saturating around $r=16$--$32$.
\textbf{Right:} visual understanding (POPE accuracy); the dashed line is the
base model. The adapter tracks the base across the whole range, i.e.\ no
cross-modal interference at any rank.}
\label{fig:rank_sweep_panels}
\end{figure}

\begin{figure}[t]
\centering
\includegraphics[width=0.7\columnwidth]{figs/kl_sweep/kl_sweep_image.pdf}
\caption{Style acquisition (CLIP directional similarity, absolute scale)
versus the KL weight $\lambda$ at rank $32$. The dotted line at $0$ marks an
adapter with no directional effect. Acquisition peaks at
$\lambda\!\approx\!0.25$--$0.5$ and collapses toward $0$ for $\lambda\!\geq\!1$,
where the adapter is over-regularised toward the base.}
\label{fig:kl_curve}
\end{figure}

\begin{figure*}[p]
\centering
\setlength{\tabcolsep}{1.5pt}\small
\begin{tabular}{cccccc}
\textbf{Base} & \textbf{$\lambda{=}0$} & \textbf{$\lambda{=}0.25$} & \textbf{$\lambda{=}0.5$} & \textbf{$\lambda{=}1.0$} & \textbf{$\lambda{=}1.5$} \\
\midrule
\includegraphics[width=0.155\textwidth]{figs/kl_comparison/base/eval_00_natural.png} &
\includegraphics[width=0.155\textwidth]{figs/kl_comparison/kl0/eval_00_natural.png} &
\includegraphics[width=0.155\textwidth]{figs/kl_comparison/kl025/eval_00_natural.png} &
\includegraphics[width=0.155\textwidth]{figs/kl_comparison/kl05/eval_00_natural.png} &
\includegraphics[width=0.155\textwidth]{figs/kl_comparison/kl1/eval_00_natural.png} &
\includegraphics[width=0.155\textwidth]{figs/kl_comparison/kl15/eval_00_natural.png} \\
\multicolumn{6}{l}{\footnotesize\textit{Prompt:} a red panda with a black face and brown eyes} \\
\addlinespace[2pt]
\includegraphics[width=0.155\textwidth]{figs/kl_comparison/base/eval_01_natural.png} &
\includegraphics[width=0.155\textwidth]{figs/kl_comparison/kl0/eval_01_natural.png} &
\includegraphics[width=0.155\textwidth]{figs/kl_comparison/kl025/eval_01_natural.png} &
\includegraphics[width=0.155\textwidth]{figs/kl_comparison/kl05/eval_01_natural.png} &
\includegraphics[width=0.155\textwidth]{figs/kl_comparison/kl1/eval_01_natural.png} &
\includegraphics[width=0.155\textwidth]{figs/kl_comparison/kl15/eval_01_natural.png} \\
\multicolumn{6}{l}{\footnotesize\textit{Prompt:} a man with glasses and a hoodie} \\
\addlinespace[2pt]
\includegraphics[width=0.155\textwidth]{figs/kl_comparison/base/human_05_natural.png} &
\includegraphics[width=0.155\textwidth]{figs/kl_comparison/kl0/human_05_natural.png} &
\includegraphics[width=0.155\textwidth]{figs/kl_comparison/kl025/human_05_natural.png} &
\includegraphics[width=0.155\textwidth]{figs/kl_comparison/kl05/human_05_natural.png} &
\includegraphics[width=0.155\textwidth]{figs/kl_comparison/kl1/human_05_natural.png} &
\includegraphics[width=0.155\textwidth]{figs/kl_comparison/kl15/human_05_natural.png} \\
\multicolumn{6}{l}{\footnotesize\textit{Prompt:} a young man with blonde spiky hair} \\
\addlinespace[2pt]
\includegraphics[width=0.155\textwidth]{figs/kl_comparison/base/human_07_natural.png} &
\includegraphics[width=0.155\textwidth]{figs/kl_comparison/kl0/human_07_natural.png} &
\includegraphics[width=0.155\textwidth]{figs/kl_comparison/kl025/human_07_natural.png} &
\includegraphics[width=0.155\textwidth]{figs/kl_comparison/kl05/human_07_natural.png} &
\includegraphics[width=0.155\textwidth]{figs/kl_comparison/kl1/human_07_natural.png} &
\includegraphics[width=0.155\textwidth]{figs/kl_comparison/kl15/human_07_natural.png} \\
\multicolumn{6}{l}{\footnotesize\textit{Prompt:} an old wizard with a long white beard} \\
\addlinespace[2pt]
\includegraphics[width=0.155\textwidth]{figs/kl_comparison/base/human_08_natural.png} &
\includegraphics[width=0.155\textwidth]{figs/kl_comparison/kl0/human_08_natural.png} &
\includegraphics[width=0.155\textwidth]{figs/kl_comparison/kl025/human_08_natural.png} &
\includegraphics[width=0.155\textwidth]{figs/kl_comparison/kl05/human_08_natural.png} &
\includegraphics[width=0.155\textwidth]{figs/kl_comparison/kl1/human_08_natural.png} &
\includegraphics[width=0.155\textwidth]{figs/kl_comparison/kl15/human_08_natural.png} \\
\multicolumn{6}{l}{\footnotesize\textit{Prompt:} a teenage girl with pink hair} \\
\end{tabular}
\caption{Effect of the KL-to-base weight $\lambda$ on visual style acquisition
(rank $32$). Style transfer is strong for $\lambda\!\in\![0,0.5]$ and visibly
collapses back toward the base for $\lambda\!\geq\!1$, where the adapter is
over-regularised.}
\label{fig:kl_visual}
\end{figure*}


\section{Text-direction rank sweep}
\label{app:textsweep}
Table~\ref{tab:textrank} sweeps the LoRA rank on the poetry target at
$\lambda{=}0.5$, the plasticity axis complementary to the anchor sweep of
Table~\ref{tab:textkl}. Increasing the rank from $4$ to $64$ steadily improves
acquisition (poetry perplexity $32.7$ to $23.8$) while MMLU, POPE and MME stay at
the base level at every rank, and WikiText perplexity is not degraded and is in
fact slightly below the base throughout. Specializing the text direction
therefore leaves the other capabilities intact across capacity, mirroring the
image result of Table~\ref{tab:rank}. Pushing the rank further to $128$ confirms
diminishing returns: acquisition improves only marginally (poetry perplexity
$23.8$ to $23.2$), while a first mild understanding cost appears (MME $1122$ to
$1105$, POPE $0.877$ to $0.867$); this is the second of the three convergent
settings of Section~\ref{sec:boundary}. We therefore report up to rank $64$,
which sits at the knee of this trade-off.

\begin{table}[t]
\centering\small
\setlength{\tabcolsep}{7pt}
\begin{tabular}{lcccccc}
\toprule
& Base & $r{=}4$ & $r{=}8$ & $r{=}16$ & $r{=}32$ & $r{=}64$ \\
\midrule
\multicolumn{7}{l}{\textit{Acquisition (target style)}}\\
Poetry PPL $\downarrow$ & $32.7$ & $25.6$ & $25.3$ & $24.9$ & $24.4$ & $\mathbf{23.8}$ \\
\midrule
\multicolumn{7}{l}{\textit{Untrained capabilities}}\\
WikiText PPL $\downarrow$ \tiny($\pm.21$) & $13.71$ & $12.52$ & $\underline{12.51}$ & $\underline{12.51}$ & $12.53$ & $12.59$ \\
MMLU $\uparrow$ \tiny($\pm.021$)          & $.481$ & $.481$ & $.477$ & $.479$ & $.477$ & $.488$ \\
POPE acc $\uparrow$ \tiny($\pm.010$)      & $.875$ & $.871$ & $.871$ & $.871$ & $.873$ & $.877$ \\
MME $\uparrow$ \tiny($\pm42$)             & $1130$ & $1129$ & $1111$ & $1128$ & $1130$ & $1122$ \\
\bottomrule
\end{tabular}
\caption{Text (poetry) rank sweep at $\lambda{=}0.5$. The style is acquired (bold
= best); knowledge and understanding stay flat. WikiText perplexity is not
degraded and is in fact slightly lower at every rank (underlined best).
Subsampled probe configuration, single seed.}
\label{tab:textrank}
\end{table}


\section{Image generation under the text adapter}
\label{app:t2itext}
\begin{figure*}[t]
\centering
\setlength{\tabcolsep}{1.5pt}
\renewcommand{\arraystretch}{0.9}
\small
\begin{tabular}{cccccc}
\textbf{Base} & $\lambda{=}0$ & $\lambda{=}0.25$ & $\lambda{=}0.5$ & $\lambda{=}1$ & $\lambda{=}1.5$ \\
\midrule
\includegraphics[width=0.15\textwidth]{figs/t2i_text/kl0/prompt0_base.png} &
\includegraphics[width=0.15\textwidth]{figs/t2i_text/kl0/prompt0_adapter.png} &
\includegraphics[width=0.15\textwidth]{figs/t2i_text/kl025/prompt0_adapter.png} &
\includegraphics[width=0.15\textwidth]{figs/t2i_text/kl05/prompt0_adapter.png} &
\includegraphics[width=0.15\textwidth]{figs/t2i_text/kl1/prompt0_adapter.png} &
\includegraphics[width=0.15\textwidth]{figs/t2i_text/kl15/prompt0_adapter.png} \\
\multicolumn{6}{l}{\footnotesize\textit{Prompt:} a cat sitting on a chair} \\
\addlinespace[2pt]
\includegraphics[width=0.15\textwidth]{figs/t2i_text/kl0/prompt1_base.png} &
\includegraphics[width=0.15\textwidth]{figs/t2i_text/kl0/prompt1_adapter.png} &
\includegraphics[width=0.15\textwidth]{figs/t2i_text/kl025/prompt1_adapter.png} &
\includegraphics[width=0.15\textwidth]{figs/t2i_text/kl05/prompt1_adapter.png} &
\includegraphics[width=0.15\textwidth]{figs/t2i_text/kl1/prompt1_adapter.png} &
\includegraphics[width=0.15\textwidth]{figs/t2i_text/kl15/prompt1_adapter.png} \\
\multicolumn{6}{l}{\footnotesize\textit{Prompt:} a medieval castle on a hill} \\
\addlinespace[2pt]
\includegraphics[width=0.15\textwidth]{figs/t2i_text/kl0/prompt2_base.png} &
\includegraphics[width=0.15\textwidth]{figs/t2i_text/kl0/prompt2_adapter.png} &
\includegraphics[width=0.15\textwidth]{figs/t2i_text/kl025/prompt2_adapter.png} &
\includegraphics[width=0.15\textwidth]{figs/t2i_text/kl05/prompt2_adapter.png} &
\includegraphics[width=0.15\textwidth]{figs/t2i_text/kl1/prompt2_adapter.png} &
\includegraphics[width=0.15\textwidth]{figs/t2i_text/kl15/prompt2_adapter.png} \\
\multicolumn{6}{l}{\footnotesize\textit{Prompt:} a bowl of fruit on a table} \\
\addlinespace[2pt]
\includegraphics[width=0.15\textwidth]{figs/t2i_text/kl0/prompt3_base.png} &
\includegraphics[width=0.15\textwidth]{figs/t2i_text/kl0/prompt3_adapter.png} &
\includegraphics[width=0.15\textwidth]{figs/t2i_text/kl025/prompt3_adapter.png} &
\includegraphics[width=0.15\textwidth]{figs/t2i_text/kl05/prompt3_adapter.png} &
\includegraphics[width=0.15\textwidth]{figs/t2i_text/kl1/prompt3_adapter.png} &
\includegraphics[width=0.15\textwidth]{figs/t2i_text/kl15/prompt3_adapter.png} \\
\multicolumn{6}{l}{\footnotesize\textit{Prompt:} an astronaut riding a horse} \\
\end{tabular}
\caption{Image generation under the text (poetry) adapter across KL weights
(rank 16). Column 1 is the base model; columns 2 to 6 are the poetry adapter at
increasing $\lambda$, on the same prompts and seeds. The generated images stay
close to the base across the whole range, the qualitative counterpart of the
near-zero Gen-CMMD in Table~\ref{tab:textkl}.}
\label{fig:t2i_text_kl}
\end{figure*}


\section{Size of the update and its weight-space localization}
\label{app:dw}
Tables~\ref{tab:dw-rank} and~\ref{tab:dw-kl} give the aggregate update sizes
quoted in \S\ref{sec:dw}. The base norm $\lVert W\rVert_F$ is read from the frozen weights;
$\lVert\Delta W\rVert_F=\lVert(\alpha/r)BA\rVert_F$ is computed from the LoRA
factors via the $r\times r$ trace identity, so the full $\Delta W$ is never
formed. Figure~\ref{fig:dw-heatmap} resolves the same quantity per module, which
is what makes the ``localized'' claim concrete rather than aggregate.

\begin{table}[h]\centering\small
\caption{Update size per adapter, $\lVert\Delta W\rVert_F/\lVert W\rVert_F$ (\%,
summed over all backbone projections), versus LoRA rank. It grows smoothly with
rank in both directions; at matched rank the image adapter changes the shared
backbone more than text.}
\label{tab:dw-rank}
\begin{tabular}{lcccccc}
\toprule
Rank $r$ & 2 & 4 & 8 & 16 & 32 & 64 \\
\midrule
Image (gen.)  & 0.58 & 0.75 & 0.97 & 1.25 & 1.57 & 1.91 \\
Text (poetry) & ---  & 0.67 & 0.84 & 1.05 & 1.32 & 1.66 \\
\bottomrule
\end{tabular}
\end{table}

\begin{table}[h]\centering\small
\caption{Update size versus the KL weight. A larger $\lambda$ shrinks the update,
more sharply for text.}
\label{tab:dw-kl}
\begin{tabular}{lccccc}
\toprule
$\lambda$ & 0 & 0.25 & 0.5 & 1 & 1.5 \\
\midrule
Image ($r{=}32$) & 1.55 & 1.58 & 1.57 & 1.45 & 1.33 \\
Text ($r{=}16$)  & 1.15 & 1.15 & 1.05 & 0.89 & 0.80 \\
\bottomrule
\end{tabular}
\end{table}

\section{The general-text perplexity anomaly}
\label{app:anomaly}
General-text perplexity falls consistently below the base after text adaptation,
across ranks, KL weights and all three targets (e.g.\ $13.71\to12.07$ for Welsh),
and the direction is corroborated by the higher MAUVE of the adapters' free
English (Table~\ref{tab:mauve}). In continual-learning terms this would be
positive backward transfer, but we have no confirmed mechanism. Plausible and
unverified causes include a mild regularization of the shared language head by
the additional text objective, and a reduced tendency to slip into the assistant
persona under free continuation, which the qualitative samples make visible: the
base frequently interrupts a continuation with ``Hello! I'm here to help''
(Appendix~\ref{app:medicine}), a behaviour the corpus-trained adapters suppress.
Separating a genuine capability gain from the removal of a persona artefact
requires multi-seed runs and a persona-controlled evaluation, neither of which we
have, so we build no claim on this observation.

\begin{figure*}[t]
\centering
\begin{subfigure}{0.32\linewidth}\centering
  \includegraphics[width=\linewidth]{figs/weight_delta/dw_heatmap_img_r16.png}
  \caption{Image gen., $r{=}16$}
\end{subfigure}\hfill
\begin{subfigure}{0.32\linewidth}\centering
  \includegraphics[width=\linewidth]{figs/weight_delta/dw_heatmap_text_r16.png}
  \caption{Text/poetry, $r{=}16$}
\end{subfigure}\hfill
\begin{subfigure}{0.32\linewidth}\centering
  \includegraphics[width=\linewidth]{figs/weight_delta/dw_heatmap_text_r64.png}
  \caption{Text/poetry, $r{=}64$}
\end{subfigure}
\caption{Per-module relative weight change
$\lVert\Delta W\rVert_F/\lVert W\rVert_F$ across layers (rows) and the seven
projections (columns). Common to both directions, the update spares the
attention query/key/value projections almost everywhere, awakening only in the
last few layers, and concentrates in the MLP. The directions differ in the fine
locus: text spreads the change broadly over \texttt{gate} across all depths,
whereas image concentrates it toward the output, peaking at the final layer's
\texttt{down} projection. Raising the rank (text $r{=}16$ vs.\ $r{=}64$) scales
the same pattern up rather than relocating it.}
\label{fig:dw-heatmap}
\end{figure*}


\section{Regularizer ablation}
\label{app:reg}
We compare the KL-to-base regularizer at fixed rank 32 in the image-generation
direction against three alternatives: L2-SP \citep{xuhong2018l2sp} (an L2 penalty
on the LoRA update, weight space), logit-MSE (a squared-error match to the base
logits, output space), and the reconstruction-loss objective of Anole
\citep{chern2024anole}, the standard recipe for fine-tuning native generators
(\texttt{CE}$+0.5\,$recon, no base anchor). All share the same training budget and
data and are evaluated identically on $80$ held-out prompts under the canonical,
trigger-free protocol used in the rest of the paper; we also report each adapter's
intervention strength $\lVert\Delta W\rVert_F/\lVert W\rVert_F$. The KL attains the
best distributional acquisition of all methods: at essentially the same
intervention strength as the unregularized baseline ($1.57$ vs.\ $1.55\%$) it
reaches a markedly lower CMMD ($.108$ vs.\ $.155$), i.e.\ it \emph{redirects} the
update rather than merely shrinking it. Logit-MSE over-constrains the output
distribution: it keeps image quality high but barely transfers the style without a
prompt trigger (DirSim $.029$, the worst of all methods), which the prompt-triggered
evaluation had masked. Weight-space L2-SP simply shrinks the update (to $0.60\%$).
The Anole reconstruction loss acquires the style comparably to the unregularized
baseline (DirSim $.162$) but with worse distributional fidelity than the KL (CMMD
$.148$ vs.\ $.108$) and, being an acquisition term rather than a base anchor, offers
no forgetting protection. The KL alone attains the best acquisition while keeping
the adapter anchored to the base. Single-seed run.

\begin{table}[t]\centering\small
\setlength{\tabcolsep}{3.5pt}
\begin{tabular}{lcccccc}
\toprule
Regularizer & $\lVert\Delta W\rVert/\lVert W\rVert$ & CMMD $\downarrow$ & KID $\downarrow$ & DINO $\uparrow$ & CLIP-IQA $\uparrow$ & DirSim $\uparrow$ \\
\midrule
None ($\lambda{=}0$) & $1.55\%$ & $.155_{\pm.009}$ & $.027_{\pm.004}$ & $.626_{\pm.013}$ & $.666_{\pm.011}$ & $.157_{\pm.007}$ \\
L2-SP                & $0.60\%$ & $.139_{\pm.009}$ & $.018_{\pm.004}$ & $.624_{\pm.011}$ & $.683_{\pm.008}$ & $\mathbf{.164}_{\pm.007}$ \\
Logit-MSE            & $1.14\%$ & $.216_{\pm.012}$ & $.106_{\pm.011}$ & $.410_{\pm.018}$ & $\mathbf{.725}_{\pm.014}$ & $.029_{\pm.006}$ \\
KL-to-base (ours)    & $1.57\%$ & $\mathbf{.108}_{\pm.010}$ & $\mathbf{.014}_{\pm.004}$ & $\mathbf{.643}_{\pm.013}$ & $.691_{\pm.010}$ & $.153_{\pm.007}$ \\
Anole (recon)       & $1.56\%$ & $.148_{\pm.009}$ & $.026_{\pm.005}$ & $.637_{\pm.011}$ & $.666_{\pm.010}$ & $.162_{\pm.007}$ \\
\bottomrule
\end{tabular}
\caption{Regularizer ablation at rank 32 (image direction, $N{=}80$, single seed;
$\pm$ one standard error; bold = best per metric). Under the canonical,
trigger-free evaluation, the KL attains the best distributional acquisition
(CMMD, KID, DINO), separating clearly from the no-regularizer baseline at matched
intervention strength ($\lVert\Delta W\rVert/\lVert W\rVert$ $1.57$ vs.\ $1.55\%$,
CMMD $.108$ vs.\ $.155$). Logit-MSE over-constrains the output distribution and
barely transfers the style without the prompt trigger (DirSim $.029$) despite high
image quality; L2-SP merely shrinks the update. The Anole reconstruction loss
acquires the style comparably to the unregularized baseline but with worse
distributional fidelity than the KL and no base anchoring.}
\label{tab:reg-ablation}
\end{table}

\clearpage
\section{Metric debugging: is the $\lambda$ optimum a bug, memorization, or noise?}
\label{app:metricdebug}
The style metrics put the image-generation acquisition optimum at $\lambda{=}0.25$,
while the eye mildly prefers $\lambda{=}0.5$. We resolve this directly by
re-scoring the \emph{already-generated} KL-sweep images (no retraining) against
both held-out and training references, and with two additional embedding metrics
(Table~\ref{tab:metric-debug}). \textbf{Not noise:} the $\lambda$-ranking of CMMD
is identical against held-out and training reals (both order
$0.25\!<\!0.5\!<\!0\!<\!1\!<\!1.5$), so the optimum is stable, not a small-$N$
artefact. \textbf{Not memorization or leakage:} for every $\lambda$ the
nearest-neighbour DINOv2 similarity to \emph{training} images is \emph{lower} than
to held-out images (and CMMD to training is higher), so the generations are not
near-copies of the training set. \textbf{A metric-choice effect, not a bug:} CMMD
(a CLIP-space distribution distance) favours $\lambda{=}0.25$, but the
memorization-robust DINOv2-space distribution distance (DINO-MMD) instead favours
$\lambda{=}0.5$ ($.035$ vs.\ $.041$), and content preservation -- the DINOv2
similarity between the adapter's image and the base image for the same prompt --
rises from $.32$ at $\lambda{=}0.25$ to $.39$ at $\lambda{=}0.5$. The CLIP style
metrics reward pure style proximity (optimum $\lambda{=}0.25$); the DINO
distribution distance and content preservation reward style transfer that keeps
the prompt content (optimum $\lambda{=}0.5$), which is what the eye weighs, and
what an independent VLM judge confirms (Table~\ref{tab:judge}). This justifies the
$\lambda{=}0.5$ Final recipe.

\begin{table}[h]
\centering\small
\setlength{\tabcolsep}{5pt}
\begin{tabular}{lcccccc}
\toprule
$\lambda$ & CMMD$_{\text{test}}\downarrow$ & CMMD$_{\text{train}}\downarrow$ & DINO-MMD$\downarrow$ & NN$_{\text{test}}\uparrow$ & NN$_{\text{train}}\uparrow$ & CtntPres$\uparrow$ \\
\midrule
$0$    & $.155$ & $.163$ & $.044$ & $.626$ & $.609$ & $.303$ \\
$0.25$ & $\mathbf{.100}$ & $.111$ & $.041$ & $.647$ & $.634$ & $.317$ \\
$0.5$  & $.108$ & $.124$ & $\mathbf{.035}$ & $.643$ & $.633$ & $.389$ \\
$1$    & $.205$ & $.221$ & $.091$ & $.439$ & $.411$ & $.721$ \\
$1.5$  & $.221$ & $.235$ & $.114$ & $.388$ & $.364$ & $.792$ \\
\bottomrule
\end{tabular}
\caption{Metric debugging on the KL sweep (rank 32), re-scoring the
already-generated images (no retraining). CMMD$_{\text{test/train}}$: CLIP-space
distribution distance to held-out vs.\ training reals (identical ranking $\Rightarrow$
signal not noise; train $>$ test $\Rightarrow$ no training over-fit). NN: nearest
DINOv2 similarity to held-out vs.\ training reals (train $<$ test everywhere
$\Rightarrow$ no memorization). DINO-MMD: memorization-robust DINOv2-space
distribution distance -- optimum $\lambda{=}0.5$. CtntPres: DINOv2 similarity
between the adapter and base image for the same prompt -- rises with $\lambda$. The
CLIP metrics favour $\lambda{=}0.25$; DINO-MMD and content preservation favour
$\lambda{=}0.5$.}
\label{tab:metric-debug}
\end{table}

As an independent corroboration we add an open VLM judge (Qwen2-VL-7B, not Janus,
to avoid self-evaluation), comparing $\lambda{=}0.25$ vs.\ $\lambda{=}0.5$ pairwise
on the same $80$ prompts in both image orders and counting only order-consistent
preferences (MT-Bench style). It prefers $\lambda{=}0.5$ on every criterion,
significantly on content (Table~\ref{tab:judge}).

\begin{table}[t]
\centering\small
\setlength{\tabcolsep}{8pt}
\begin{tabular}{lccc}
\toprule
Criterion & $\lambda{=}0.25$ win & $\lambda{=}0.5$ win & tie \\
\midrule
Style    & $23.8\%$ & $\mathbf{30.0\%}$ & $46.2\%$ \\
Content  & $6.2\%$  & $\mathbf{20.0\%}$ & $73.8\%$ \\
Overall  & $23.8\%$ & $\mathbf{30.0\%}$ & $46.2\%$ \\
\bottomrule
\end{tabular}
\caption{Independent VLM judge (Qwen2-VL-7B), $\lambda{=}0.25$ vs.\ $\lambda{=}0.5$
on $80$ held-out prompts, each scored in both image orders. A configuration wins a
prompt only if preferred in \emph{both} orders; order-dependent results are ties
(position-bias mitigation). On the order-consistent comparisons the judge prefers
$\lambda{=}0.5$ on every criterion, significantly on content (decisive win-rate
$76\%$, $95\%$ CI $[55,89]$), and never favours $\lambda{=}0.25$.}
\label{tab:judge}
\end{table}

\clearpage
\section{Text-direction qualitative examples}
\label{app:textqual}
We show base-vs-adapter continuations for the poetry text-LoRA ($r{=}16$,
$\lambda{=}0.5$); the base is the identical model with the adapter disabled, so the
only change between columns is the LoRA. Decoding is greedy with a repetition
penalty of $1.3$ (Appendix~\ref{app:impl}). Each incipit is continued from a fixed
two-line verse couplet given as hidden context to \emph{both} models (so the
unified model stays in a poem rather than its assistant persona); we display only
the text generated after the incipit. With the same primed context the adapter
sustains metrical, lineated verse, whereas the base drifts into prose or stops
short. On bare, non-poetic prompts the register does not surface: the LoRA is a
next-token language model, not instruction-tuned.

\begin{table}[t]\centering\small
\renewcommand{\arraystretch}{1.3}
\begin{tabular}{@{}p{0.45\linewidth}p{0.45\linewidth}@{}}
\toprule
\multicolumn{2}{@{}l}{\textit{Incipit:} ``The autumn wind\,\dots''}\\
\midrule
\textbf{Base (adapter off)} & \textbf{Poetry LoRA}\\[2pt]
\footnotesize is blowing in our faces; it's a time to reflect on what has been
accomplished this year\,\dots
&
\footnotesize is blowing in our faces; it's a-blowing! ,  O what shall I do?\\
\midrule
\multicolumn{2}{@{}l}{\textit{Incipit:} ``Upon the silent hill there\,\dots''}\\
\midrule
\textbf{Base (adapter off)} & \textbf{Poetry LoRA}\\[2pt]
\footnotesize lies a light of gold; ,  O! how it gleams in twilight's shade: it
is not day nor night but something more divine.
&
\footnotesize stands a tree;\newline
The leaves are green as emeralds in their sleep:\newline
And on its branches hang an empty cage,\newline
A prison for some bird of sorrow's keep!\\
\midrule
\multicolumn{2}{@{}l}{\textit{Incipit:} ``Beneath the ancient oak,\,\dots''}\\
\midrule
\textbf{Base (adapter off)} & \textbf{Poetry LoRA}\\[2pt]
\footnotesize where time has worn its way; --- O how I love thee!
&
\footnotesize where I have stood;\newline
The sun was setting on a summer's day: ---\newline
I saw her face for one brief moment flash!\newline
And then she vanished from before mine eyes.\\
\midrule
\multicolumn{2}{@{}l}{\textit{Incipit:} ``O gentle heart, what\,\dots''}\\
\midrule
\textbf{Base (adapter off)} & \textbf{Poetry LoRA}\\[2pt]
\footnotesize words of love can say? --- O! for a thousand tongues to tell thee how
I feel this day; but none are fit nor true \dots\ (1 Samuel 25).
&
\footnotesize is thy secret? What art thou made of but a little child's delight in
play; \dots\ O say not ``I am nothing,'' for I see how much you have to give: your
eyes shine like stars on earth when they look at me\,\dots\\
\bottomrule
\end{tabular}
\caption{Text-direction qualitative examples (poetry LoRA, $r{=}16$,
$\lambda{=}0.5$; greedy decoding, repetition penalty $1.3$). Each incipit is
continued from a fixed verse couplet given as hidden context to \emph{both}
models; we show only the text after the incipit. With the same primed context the
adapter sustains lineated, imagistic verse (e.g.\ the emerald-leaf quatrain),
whereas the base is shorter, more abstract, or drifts into prose (top) or
scriptural citation (bottom).}
\label{tab:text-examples}
\end{table}


\section{Medicine: qualitative before and after}
\label{app:medicine}
Same prompts, same model, adapter off (base) vs.\ on (medicine LoRA). The text
continuations show the register shift toward clinical prose; the image grid shows
that text specialisation leaves the generation direction essentially unchanged.

\begin{table}[t]
\centering\footnotesize
\setlength{\tabcolsep}{6pt}\renewcommand{\arraystretch}{1.25}
\begin{tabular}{p{0.45\textwidth} p{0.45\textwidth}}
\toprule
\textbf{Base} & \textbf{Medicine LoRA} \\
\midrule
\multicolumn{2}{@{}p{0.95\textwidth}@{}}{\textit{Prompt:} ``What is the mechanism of action of metformin?''}\\
Metformin works by inhibiting glucose production in cells, particularly those
that are most active at high concentrations. Hello! I'm here to help you with
your question [\dots]
&
Metformin works by inhibiting glucose uptake in muscle cells and reducing
hepatic gluconeogenesis. [\dots] \\
\midrule
\multicolumn{2}{@{}p{0.95\textwidth}@{}}{\textit{Prompt:} ``A 58-year-old man with hypertension presents with acute chest pain radiating to the left arm; the most likely diagnosis is\dots''}\\
\textit{(returns another question instead of a diagnosis:)} ``What are some
common causes of sudden cardiac death in older adults?''
&
unstable angina. A coronary catheterization should be performed as soon as
possible and an appropriate treatment plan must begin immediately; if there are
no signs of myocardial infarction (MI), an electrocardiogram (ECG) and stress
test [\dots] \\
\midrule
\multicolumn{2}{@{}p{0.95\textwidth}@{}}{\textit{Prompt (general-knowledge control):} ``Who wrote the play Romeo and Juliet?''}\\
William Shakespeare. [\dots] the playwright of this famous tragedy [\dots]
($1564$--$1623$).
&
William Shakespeare. [\dots] a famous English playwright who lived between
approximately $1564$--$1616$ [\dots] including Romeo \& Juliet! \\
\bottomrule
\end{tabular}
\caption{Text continuations (greedy, deterministic; lightly trimmed, [\dots]).
On medical prompts the base stays vague or fails to answer, while the medicine
adapter gives precise clinical content and the register of PubMed (drug
mechanism, ``reducing hepatic gluconeogenesis''; a diagnosis with management). On
a general-knowledge control both answer correctly, i.e.\ the domain adaptation
does not cost general factual recall. Residual chat artifacts (``Hello!\dots'')
appear because the base is partly instruction-tuned.}
\label{tab:med-text}
\end{table}

\begin{figure}[t]
\centering
\setlength{\tabcolsep}{1.5pt}\renewcommand{\arraystretch}{1.0}
\begin{tabular}{@{}cccc@{}}
\rotatebox{90}{\footnotesize\ \ Base} &
\includegraphics[width=0.30\textwidth]{figs/med_comparison/prompt0_base.png} &
\includegraphics[width=0.30\textwidth]{figs/med_comparison/prompt1_base.png} &
\includegraphics[width=0.30\textwidth]{figs/med_comparison/prompt6_base.png} \\
\rotatebox{90}{\footnotesize Med LoRA} &
\includegraphics[width=0.30\textwidth]{figs/med_comparison/prompt0_adapter.png} &
\includegraphics[width=0.30\textwidth]{figs/med_comparison/prompt1_adapter.png} &
\includegraphics[width=0.30\textwidth]{figs/med_comparison/prompt6_adapter.png} \\
\end{tabular}
\caption{Image generation under the text--medicine LoRA: same prompts and seeds,
adapter off (top) vs.\ on (bottom). Generations are essentially unchanged,
confirming the text adapter does not degrade the generation direction (cf.\ the
near-zero Gen-CMMD).}
\label{fig:med-t2i}
\end{figure}


\section{Welsh text-direction qualitative examples}
\label{app:welsh}
Base-vs-adapter continuations for the Welsh text-LoRA ($r{=}32$, $\lambda{=}0.1$);
the base is the identical model with the adapter disabled. Greedy decoding,
repetition penalty $1.3$. Two readings: (i) on \emph{Welsh} seeds the adapter
sustains Welsh while the base breaks into English/assistant text (acquisition);
(ii) on the \emph{English} seed the adapter stays in English (no bleed). As with
poetry, greedy decoding is locally repetitive; sampled decoding ($T{=}0.9$, last
row) is more fluent.

\begin{table}[t]\centering\small
\renewcommand{\arraystretch}{1.3}
\begin{tabular}{@{}p{0.45\linewidth}p{0.45\linewidth}@{}}
\toprule
\multicolumn{2}{@{}l}{\textit{Welsh seed:} ``Mae Cymru yn wlad\,\dots'' (\textit{Wales is a country})}\\
\midrule
\textbf{Base (adapter off)} & \textbf{Welsh LoRA}\\[2pt]
\footnotesize fod â'r eich, a dda i gael \dots\ Hello Mae! I am happy to help you
with your Welsh language question.
&
\footnotesize fwyaf o'r gogledd-ddinasol a ddau 10,578.234 iawn (a chyfansoddwyd
gan yr Aelodau Canolfan Llyfrau Genedlaethol). Mae hynny wedi ei bod arall am\,\dots\\
\midrule
\multicolumn{2}{@{}l}{\textit{English seed:} ``The history of Wales\,\dots''}\\
\midrule
\textbf{Base (adapter off)} & \textbf{Welsh LoRA}\\[2pt]
\footnotesize is a long and complex one, with many different cultures influencing
its development over the centuries. Hello!
&
\footnotesize is a long and complex one, with many different cultures influencing
the country over time. I will be writing about some key events in Welsh History
that have shaped our nation today\,\dots\\
\midrule
\multicolumn{2}{@{}l}{\textit{Welsh seed (sampled, $T{=}0.9$):} ``Yr oedd unwaith\,\dots'' (\textit{There was once})}\\
\midrule
\textbf{Base (adapter off)} & \textbf{Welsh LoRA}\\[2pt]
\footnotesize Hello world \dots\ Wyt yn agor fod dinio? Nochthawd, bye!
&
\footnotesize a byddai'n anhysbys, neu llawn barwas yng ngogleddol rhagolig sydd
wedi bod yn enwu am fod fel cysgu gerfansau pencampell\,\dots\\
\bottomrule
\end{tabular}
\caption{Welsh text-direction qualitative examples ($r{=}32$, $\lambda{=}0.1$). On
Welsh seeds the adapter sustains Welsh where the base slips into English/assistant
text; on the English seed it stays in English (no bleed). Greedy rows are locally
repetitive; the sampled row is more fluent.}
\label{tab:welsh-examples}
\end{table}

\begin{figure}[t]\centering
\setlength{\tabcolsep}{1.5pt}\renewcommand{\arraystretch}{0.8}\small
\begin{tabular}{cc}
\textbf{Base} & \textbf{Welsh LoRA}\\
\includegraphics[width=0.30\linewidth]{figs/t2i_text/welsh/prompt0_base.png} &
\includegraphics[width=0.30\linewidth]{figs/t2i_text/welsh/prompt0_adapter.png}\\
\multicolumn{2}{c}{\scriptsize\textit{a red fox sitting in a snowy forest at dawn}}\\
\addlinespace[1pt]
\includegraphics[width=0.30\linewidth]{figs/t2i_text/welsh/prompt2_base.png} &
\includegraphics[width=0.30\linewidth]{figs/t2i_text/welsh/prompt2_adapter.png}\\
\multicolumn{2}{c}{\scriptsize\textit{a close-up portrait of an elderly fisherman with a wool cap}}\\
\addlinespace[1pt]
\includegraphics[width=0.30\linewidth]{figs/t2i_text/welsh/prompt4_base.png} &
\includegraphics[width=0.30\linewidth]{figs/t2i_text/welsh/prompt4_adapter.png}\\
\multicolumn{2}{c}{\scriptsize\textit{a futuristic city skyline at night with neon lights}}\\
\addlinespace[1pt]
\includegraphics[width=0.30\linewidth]{figs/t2i_text/welsh/prompt6_base.png} &
\includegraphics[width=0.30\linewidth]{figs/t2i_text/welsh/prompt6_adapter.png}\\
\multicolumn{2}{c}{\scriptsize\textit{a plate of spaghetti with tomato sauce and fresh basil}}\\
\addlinespace[1pt]
\includegraphics[width=0.30\linewidth]{figs/t2i_text/welsh/prompt7_base.png} &
\includegraphics[width=0.30\linewidth]{figs/t2i_text/welsh/prompt7_adapter.png}\\
\multicolumn{2}{c}{\scriptsize\textit{an astronaut riding a horse on the surface of Mars}}\\
\end{tabular}
\caption{Image generation under the Welsh text adapter ($r{=}32$, $\lambda{=}0.1$),
base vs adapter on identical prompts and seeds. Generations stay close to the base,
consistent with the near-zero Gen-CMMD ($1.7\times10^{-3}$,
Table~\ref{tab:targets}).}
\label{fig:t2i_welsh}
\end{figure}


\section{Free-form generation samples (MAUVE)}
\label{app:mauve}
The base's low absolute MAUVE-to-WikiText ($0.38$) is not degenerate text: base
continuations are coherent, encyclopedic English (each starts from a fixed
WikiText prefix, then the model continues). MAUVE values are relative to our fixed
protocol ($N{\le}2000$, GPT-2-large features) and are meaningful for the
within-table comparison, not as absolute cross-paper numbers.

\begin{table}[h]\centering\footnotesize
\renewcommand{\arraystretch}{1.2}
\begin{tabular}{@{}p{0.95\linewidth}@{}}
\toprule
\textbf{Base, free-form English continuations (WikiText prefix + generation)}\\
\midrule
\footnotesize Sails persisted on `cruising ironclads' for much longer. During the
1860s, the French navy had produced the Alma and La V\'erendaye, two of the first
ironclads to be built\dots\\
\midrule
\footnotesize Mortier accepted the bait of a rumored Russian retreat. In the early
morning of 11 November he and Gazan departed from D\"urenstein to seize the
abandoned fort\dots\\
\midrule
\footnotesize Following a hiatus of two years, Louis Neimand was hired as rabbi in
1963. Born in New York City in 1919, Neimand grew up in Brooklyn and graduated
from Cornell University in 1941\dots\\
\bottomrule
\end{tabular}
\caption{Base free-form English generations are coherent WikiText-style prose
(cf.\ MAUVE $0.38$ in Table~\ref{tab:mauve}); the low absolute value reflects the
distributional gap of a multimodal assistant model from encyclopedic text, not
degeneracy.}
\label{tab:mauve-samples}
\end{table}


\section{Cross-modal transfer of the acquired language}
\label{app:xmodal}
The probes reported in Table~\ref{tab:xmodal} are constructed as follows. For (A)
we generate the same content from parallel English/Welsh prompts
(Table~\ref{tab:xmodal-welsh-prompts}) and score every image against the same
\emph{English} reference with CLIPScore (a strong English CLIP; the Welsh side is
thus scored on meaning, not on a Welsh-blind text encoder). For (C) the image and
answer format are fixed and only the question language changes; answers are
scored by the Yes/No logit, so output-language drift cannot corrupt them. Because
the base barely knows Welsh, an improvement would be positive cross-modal
transfer of the new language, not a retention effect.

\begin{table}[h]\centering\footnotesize
\setlength{\tabcolsep}{6pt}
\begin{tabular}{@{}p{0.46\linewidth}p{0.46\linewidth}@{}}
\toprule
\textbf{English prompt} & \textbf{Welsh prompt} \\
\midrule
a red fox sitting in a snowy forest & llwynog coch yn eistedd mewn coedwig eira \\
a yellow school bus on a city street & bws ysgol melyn ar stryd y ddinas \\
a slice of chocolate cake on a plate & tafell o gacen siocled ar bl\^at \\
an old stone castle on a hill & castell carreg hen ar fryn \\
a lighthouse on a rocky coast at sunset & goleudy ar arfordir creigiog ar fachlud \\
a plate of fish and chips & pl\^at o bysgod a sglodion \\
\bottomrule
\end{tabular}
\caption{Sample of the parallel English/Welsh prompts for the text$\to$image
probe (Table~\ref{tab:xmodal}, A); the full set has $30$ pairs. The Welsh is
a translation of the English, and both are scored against the same English
reference.}
\label{tab:xmodal-welsh-prompts}
\end{table}

\begin{figure}[h]\centering
\setlength{\tabcolsep}{1.5pt}\renewcommand{\arraystretch}{0.9}\footnotesize
\begin{tabular}{ccc}
\textbf{English $\to$ Base} & \textbf{Welsh $\to$ Base} & \textbf{Welsh $\to$ Welsh LoRA} \\
\includegraphics[width=0.3\linewidth]{figs/xmodal_welsh/p00_s0_base_en.png} &
\includegraphics[width=0.3\linewidth]{figs/xmodal_welsh/p00_s0_base_cy.png} &
\includegraphics[width=0.3\linewidth]{figs/xmodal_welsh/p00_s0_adapter_cy.png} \\
\multicolumn{3}{l}{\scriptsize\textit{a red fox\dots / llwynog coch yn eistedd mewn coedwig eira}} \\
\addlinespace[2pt]
\includegraphics[width=0.3\linewidth]{figs/xmodal_welsh/p02_s0_base_en.png} &
\includegraphics[width=0.3\linewidth]{figs/xmodal_welsh/p02_s0_base_cy.png} &
\includegraphics[width=0.3\linewidth]{figs/xmodal_welsh/p02_s0_adapter_cy.png} \\
\multicolumn{3}{l}{\scriptsize\textit{a yellow school bus\dots / bws ysgol melyn ar stryd y ddinas}} \\
\addlinespace[2pt]
\includegraphics[width=0.3\linewidth]{figs/xmodal_welsh/p02_s0_base_en.png} &
\includegraphics[width=0.3\linewidth]{figs/xmodal_welsh/p02_s0_base_cy.png} &
\includegraphics[width=0.3\linewidth]{figs/xmodal_welsh/p02_s0_adapter_cy.png} \\
\multicolumn{3}{l}{\scriptsize\textit{a yellow school bus\dots / bws ysgol melyn ar stryd y ddinas}} \\
\addlinespace[2pt]
\includegraphics[width=0.3\linewidth]{figs/xmodal_welsh/p04_s0_base_en.png} &
\includegraphics[width=0.3\linewidth]{figs/xmodal_welsh/p04_s0_base_cy.png} &
\includegraphics[width=0.3\linewidth]{figs/xmodal_welsh/p04_s0_adapter_cy.png} \\
\multicolumn{3}{l}{\scriptsize\textit{a slice of chocolate cake\dots / tafell o gacen siocled ar bl\^at}} \\
\end{tabular}
\caption{Text$\to$image under the Welsh adapter. Left: the English prompt renders
correctly (base). Middle and right: the same content requested in Welsh fails for
both base and adapter, i.e.\ the acquired language does not transfer to the
generation pathway (cf.\ Table~\ref{tab:xmodal}).}
\label{fig:xmodal-welsh}
\end{figure}

\end{document}

